% concatenated from main.tex, preamble.tex
\documentclass{amsart}

% Set margin size (amsart default is about 4.5cm)
	%\usepackage[margin=1.3in]{geometry}

% Accept different input encodings
	\usepackage[utf8]{inputenc}

% AMS packages
	\usepackage{amsmath,amsfonts,amssymb,amsthm}

% Allowing bold & small caps
	\usepackage[T1]{fontenc}

% Bold section headings
	\usepackage{etoolbox}
	\patchcmd{\section}{\scshape}{\scshape\bfseries}{}{}
	\makeatletter
	\renewcommand{\@secnumfont}{\scshape\bfseries}
	\makeatother

% Empty set
	\let\oldemptyset\emptyset
	\let\emptyset\varnothing
	
% Varepsilon
	\let\oldepsilon\epsilon
	\let\epsilon\varepsilon

% Remove indentions and add space between paragraphs
	\usepackage[indent=10pt]{parskip}

% Link and bookmark setup
	\usepackage[bookmarks=true,bookmarksopen=true]{hyperref}

% Theorem style and numbering
%	\theoremstyle{theorem}
	\numberwithin{equation}{section}
    
    \newtheorem{theorem}[equation]{Theorem}    
	\newtheorem{claim}[equation]{Claim}
	\newtheorem{conjecture}[equation]{Conjecture}
	\newtheorem*{question}{Question}
	\newtheorem{corollary}[equation]{Corollary}
	\newtheorem{lemma}[equation]{Lemma}
	\newtheorem{proposition}[equation]{Proposition}
		
	\theoremstyle{plain}
	\newtheorem{result}{Theorem}
	\newtheorem{propositionresult}[result]{Proposition}
	\newtheorem{corollaryresult}[result]{Corollary}
	\renewcommand*{\theresult}{\Alph{result}}
	
	\theoremstyle{definition}
	\newtheorem*{definition}{Definition}
	\newtheorem{example}[equation]{Example}
	
	\theoremstyle{remark}
	\newtheorem{remark}[equation]{Remark}
	
%	\numberwithin{remark}{section}
	
	
	
	\numberwithin{table}{section}
	
	\renewcommand\qedsymbol{$\blacksquare$}

% Tikz setup
	\usepackage{tikz}
	\usetikzlibrary{matrix}

% Todo notes
	\usepackage{todonotes}

% Enumerate with letters
	\usepackage{enumerate}

% Table and figure setup
	\usepackage{float}

% Footnote setup
	\makeatletter
	\def\blfootnote{\gdef\@thefnmark{}\@footnotetext}
	\makeatother

% Definitions of operators
	\DeclareMathOperator{\Ad}{Ad}
	\DeclareMathOperator{\codim}{codim}
	\DeclareMathOperator{\cohom}{cohom}
	\DeclareMathOperator{\curv}{curv}
	\DeclareMathOperator{\diag}{diag}
	\DeclareMathOperator{\dist}{dist}
	\DeclareMathOperator{\grad}{grad}
	\DeclareMathOperator{\Hess}{Hess}
    \DeclareMathOperator{\Id}{Id}
	\DeclareMathOperator{\II}{I\! I}
	\DeclareMathOperator{\Image}{Im}
	\DeclareMathOperator{\Iso}{Iso}
	\DeclareMathOperator{\Isot}{Isot}
	\DeclareMathOperator{\Ker}{Ker}
	\DeclareMathOperator{\rad}{rad}
	\DeclareMathOperator{\rank}{rank}
	\DeclareMathOperator{\Ric}{Ric}
        \DeclareMathOperator{\Scal}{Scal}
	\DeclareMathOperator{\spann}{span}
	\DeclareMathOperator{\symrank}{symrank}
	\DeclareMathOperator{\vol}{vol}
	\DeclareMathOperator{\trace}{trace}

% Relations, functions, etc.
	\newcommand{\tang}{\top}
	\newcommand{\inj}{\hookrightarrow}
	\newcommand{\surj}{\twoheadrightarrow}

% Macros
	\newcommand{\of}[1]{\left( #1 \right)}
	\newcommand{\ceil}[1]{\left\lceil #1 \right\rceil}
	\newcommand{\floor}[1]{\left\lfloor #1 \right\rfloor}
	\newcommand{\pfrac}[2]{\left( \frac{#1}{#2} \right)}
	\newcommand{\inner}[1]{\left\langle #1 \right\rangle}
	\newcommand{\gen}[1]{\left\langle #1 \right\rangle}
	\newcommand{\pr}[1]{\Ric_{#1} > 0}

% Lie groups (sans serif)
	\newcommand{\Esix}{\mathsf{E}_6}
	\newcommand{\Eseven}{\mathsf{E}_7}
	\newcommand{\Eeight}{\mathsf{E}_8}
	\newcommand{\Ffour}{\mathsf{F}_4}
	\newcommand{\gG}{\mathsf{G}}
	\newcommand{\GL}{\mathsf{GL}}
	\newcommand{\gH}{\mathsf{H}}
	\newcommand{\gK}{\mathsf{K}}
	\newcommand{\gL}{\mathsf{L}}
	\newcommand{\gO}{\mathsf{O}}
	\newcommand{\PU}{\mathsf{PU}}
	\newcommand{\gS}{\mathsf{S}}
	\newcommand{\SL}{\mathsf{SL}}
	\newcommand{\SO}{\mathsf{SO}}
	\newcommand{\SOO}[2]{\text{S(O(#1)O(#2))}}
	\newcommand{\Sp}{\mathsf{Sp}}
	\newcommand{\Spin}{\mathsf{Spin}}
	\newcommand{\SU}{\mathsf{SU}}
	\newcommand{\gT}{\mathsf{T}}
	\newcommand{\gU}{\mathsf{U}}
	
% Lie algebras (fraktur)
	\newcommand{\fa}{\mathfrak{a}}
	\newcommand{\fb}{\mathfrak{b}}
	\newcommand{\ff}{\mathfrak{f}}
	\newcommand{\fg}{\mathfrak{g}}
	\newcommand{\fgl}{\mathfrak{gl}}
	\newcommand{\fh}{\mathfrak{h}}
	\newcommand{\fk}{\mathfrak{k}}
	\newcommand{\fl}{\mathfrak{l}}
	\newcommand{\fm}{\mathfrak{m}}
	\newcommand{\fn}{\mathfrak{n}}
	\newcommand{\fp}{\mathfrak{p}}
	\newcommand{\fr}{\mathfrak{r}}
	\newcommand{\so}{\mathfrak{so}}
	\renewcommand{\sp}{\mathfrak{sp}}
    \newcommand{\spin}{\mathfrak{spin}}
	\newcommand{\su}{\mathfrak{su}}
	\newcommand{\ft}{\mathfrak{t}}
	\newcommand{\fu}{\mathfrak{u}}
	\newcommand{\fz}{\mathfrak{z}}
	
% Symmetric spaces
	\newcommand{\RP}{\mathbb{R}\mathrm{P}}
	\newcommand{\CP}{\mathbb{C}\mathrm{P}}
	\newcommand{\mCP}[1]{\overline{\mathbb{C}\mathrm{P}^{#1}}}
	\newcommand{\HP}{\mathbb{H}\mathrm{P}}
	\newcommand{\OP}{\mathbb{O}\mathrm{P}^2}
	\newcommand{\Gr}{\mathrm{Gr}}

% Rings, fields, etc. (blackboard bold)
	\newcommand{\Ca}{\mathrm{Ca}}
	\newcommand{\bC}{\mathbb{C}}
	\newcommand{\bD}{\mathbb{D}}
	\newcommand{\bF}{\mathbb{F}}
	\newcommand{\bH}{\mathbb{H}}
	\newcommand{\bK}{\mathbb{K}}
	\newcommand{\bN}{\mathbb{N}}
	\newcommand{\bO}{\mathbb{O}}
	\newcommand{\bP}{\mathbb{P}}
	\newcommand{\bQ}{\mathbb{Q}}
	\newcommand{\bR}{\mathbb{R}}
	\newcommand{\bS}{\mathbb{S}}
	\newcommand{\bZ}{\mathbb{Z}}
	
% Caligraphic
	\newcommand{\cD}{\mathcal{D}}
	\newcommand{\cF}{\mathcal{F}}
	\newcommand{\cH}{\mathcal{H}}
	\newcommand{\cJ}{\mathcal{J}}
	\newcommand{\cL}{\mathcal{L}}
	\newcommand{\cV}{\mathcal{V}}
	
% Greek letters
	\newcommand{\ep}{\epsilon}
	
% Replace l in math mode with \ell
	\mathcode`l="8000
	\begingroup
	\makeatletter
	\lccode`\~=`\l
	\DeclareMathSymbol{\lsb@l}{\mathalpha}{letters}{`l}
	\lowercase{\gdef~{\ifnum\the\mathgroup=\m@ne \ell \else \lsb@l \fi}}%
	\endgroup

% Define colon equal (:=)
	\newcommand*{\defeq}{\mathrel{\vcenter{\baselineskip0.5ex \lineskiplimit0pt
				\hbox{\scriptsize.}\hbox{\scriptsize.}}}%
		=}

\usepackage[foot]{amsaddr}

\title[$\gG/\gH$ with two isotropy summands develops $\Ric>0$ under Ricci flow]{Homogeneous spaces with two equivalent isotropy summands develop positive Ricci curvature under Ricci flow}

\author[E. Cochran]{Eric Cochran$^1$}
\address{Department of Mathematics, University of California, Irvine, USA$^1$}
\email{ecochra2@uci.edu}

\author[A. Mingajev]{Arseny Mingajev$^{2}$}
\address{Department of Mathematics, Trinity University, San Antonio, TX, USA$^2$}
\email{amingaje@trinity.edu}

\author[L. Mouillé]{Lawrence Mouillé$^2$}
%\address{Trinity University, San Antonio, TX, USA}
\email{lmouille@trinity.edu}

\author[N. Valiyakath]{Nazia Valiyakath$^3$}
\address{Department of Mathematics, University of Rochester, Rochester, NY, USA$^3$}
\email{naziavaliyakath@gmail.com}

\subjclass[2020]{53E20, 53C20, 53C21, 53C30}

\date{\today}

\begin{document}

\begin{abstract}
    We study normalized Ricci flow on simply connected homogeneous spaces $\gG/\gH$ for which the isotropy representation splits into exactly two equivalent irreducible subrepresentations.
    We prove that every $\gG$-invariant metric evolves to one with positive Ricci curvature, and that the family of $\gG$-invariant metrics with positive Ricci curvature is forward-invariant under the flow. 
    The proof relies on the fact that the phase portrait of the family of fixed-volume $\gG$-invariant metrics can be explicitly visualized. 
\end{abstract}

\maketitle

\section{Introduction}

    % \todo[inline]{To post version 1 on arXiv: (1) abstract. (2) Intro. (3) Proof of $\Ric>0$ result. (3.1) move x,y coordinate change earlier (start of section 3), (3.2) update pictures, (3.3) check proof of 3.8, (3.4) have we proven/remarked that invariant curves move into P>0?, (3.5) road map of proof in Section 3 intro. (4) Make exposition flow / make notation and terminology consistent.}
    % \todo[inline]{Version 2 on arXiv: (1) ID asymptotes of system.}
    % \todo[inline]{Version 3: (1) Ledger-Obata spaces. (2) Blow downs of metrics at infinity.}

Ricci flow, introduced by Hamilton in \cite{Hamilton82}, constitutes a powerful geometric evolution equation for deforming Riemannian metrics. 
Given a Riemannian manifold $(M,g_0)$, the flow is defined by the partial differential equation
\[
    \frac{\partial}{\partial t} g(t) = -2 \Ric_{g(t)}, \qquad g(0) = g_0.
\]
Heuristically, much like the heat equation, Ricci flow is expected to have a regularizing effect, evolving a given Riemannian metric to some \textit{canonical} metric on $M$.  
A general goal is to then draw topological or geometric conclusions from the existence of such a metric.
When combined with Perelman’s monotonicity functionals and entropy techniques, Ricci flow has been instrumental in resolving several long-standing conjectures in geometry and topology, most notably the Poincaré and Geometrization Conjectures.

Ricci flow is known to preserve several important curvature conditions, including nonnegative scalar curvature \cite{Hamilton82}, positive-semidefinite curvature operator \cite{Hamilton86}, and nonnegative isotropic curvature \cite{BrendleSchoen09,Hamilton97,Nguyen10}.
Brendle and Schoen subsequently used this last property in their ground-breaking proof of the Differentiable Sphere Theorem in \cite{BrendleSchoen09}.
A unified approach to proving invariance of nonnegative curvature conditions was developed by Wilking in \cite{Wilking13}
% Ricci flow is known to preserve several important curvature conditions.
% Hamilton showed that it preserves nonnegative scalar curvature in \cite{Hamilton82} and positive-semidefinite curvature operator in \cite{Hamilton86}.
% He also showed it preserves nonnegative isotropic curvature for dimension four in \cite{Hamilton97}, which was extended to all dimensions independently by Brendle and Schoen in \cite{BrendleSchoen09} and Nguyen in \cite{Nguyen10}.
% Brendle and Schoen subsequently used this property in their ground-breaking proof of the Differentiable Sphere Theorem.
% A unified approach to proving invariance of nonnegative curvature conditions was developed by Wilking in \cite{Wilking13}
However, not all lower curvature bounds are preserved under the flow; see the work of Ni \cite{Ni04}, B\"ohm and Wilking \cite{BohmWilking07}, M\'aximo \cite{Maximo11,Maximo14}, and Bettiol and Krishnan \cite{BettiolKrishnan19,BettiolKrishnan23}.
% However, not all lower curvature bounds are preserved under the flow. 
% Ni constructed examples in \cite{Ni04} of complete, non-compact manifolds with nonnegative sectional curvature whose metrics evolve under Ricci flow to exhibit mixed sectional curvatures. 
% In \cite{BohmWilking07}, B\"ohm and Wilking showed that on the compact homogeneous space $\Sp(3)/\Sp(1)\Sp(1)\Sp(1)$, Ricci flow evolves certain metrics with postive sectional curvature into metrics with mixed Ricci curvature.
% M\'aximo gave examples in \cite{Maximo11} (resp. \cite{Maximo14}) of K\"ahler manifolds with non-negative (resp. positive) Ricci curvature that lose this property under Ricci flow.
% Bettiol and Krishnan gave examples in \cite{BettiolKrishnan19} (resp. \cite{BettiolKrishnan23}) of four-dimensional spaces with non-negative (resp. positive) sectional curvature which evolve to have mixed sectional curvatures. 

An important property of Ricci flow is that it preserves the symmetries of the initial metric $g_0$.
This property follows from the fact that the Ricci tensor is equivariant (or natural) with respect to pullbacks by diffeomorphisms of the manifold. 
Consequently, on homogeneous spaces $\gG/\gH$, Ricci flow reduces to a system of ordinary differential equations when restricted to the the space of $\gG$-invariant metrics. 
In cases that the isotropy representation decomposes into few irreducible subrepresentations (isotropy summands), the space of $\gG$-invariant metrics is low-dimensional.
Furthermore, by considering normalized Ricci flow given by
\[
    \frac{\partial}{\partial t} g(t) = -2 \Ric_{g(t)} + \frac{2 \Scal}{\dim (M)} g(t), \qquad g(0)=g_0,
\]
the family of $\gG$-invariant metrics with normalized volume is preserved under the flow.
Restricting to such a family of metrics produces an even lower-dimensional system to analyze, and the dynamics of normalized Ricci flow on this entire space of metrics can be completely visualized.

This reduction to low-dimensional ODE systems has been exploited by several authors to study the preservation of curvature conditions.
Abiev and Nikonorov investigated the preservation of positive sectional and Ricci curvature on Wallach spaces in \cite{AbievNikonorov16}, 
Abiev studied Stiefel manifolds $\SO(n)/\SO(n-2)$ in \cite{Abiev25b}, and DeVito, Gonz\'alez-\'Alvaro, and Zarei focused on spheres and complex projective spaces in \cite{DGAZpreprint}.
For more results on homogeneous Ricci flow, see the work of Abiev \cite{Abiev25a}, Abiev et.~al.~\cite{AANS14}, B\"ohm \cite{Bohm15}, B\"ohm and Lafuente \cite{BohmLafuente18}, Buzano \cite{Buzano14}, Cavenaghi et.~al.~\cite{CGM}, Gonz\'alez-\'Alvaro, and Zarei \cite{GAZ24,GAZ25}, Lauret \cite{Lauret13}, Sbiti \cite{Sbiti20}, and Statha \cite{Statha22}.

% This reduction to low-dimensional ODE systems has been exploited by several authors to study the preservation of curvature conditions.
% Abiev and Nikonorov studied normalized Ricci flow on the Wallach spaces $W^6 = \SU(3) / \gT^2_{\mathrm{max}}$, $W^{12} = \Sp(3)/\Sp(1) \Sp(1) \Sp(1)$, and $W^{24} = \Ffour / \Spin(8)$, whose isotropy representations decompose into three inequivalent summands.
% They showed that generic initial metrics with positive sectional curvature will evolve to have mixed sectional curvatures, that an analogous statement holds for Ricci curvature on the $W^{12}$ and $W^{24}$, and that under certain conditions, positivity of Ricci curvature is preserved on generalized Wallach spaces.
% Abiev studied normalized Ricci flow on the Stiefel manifolds $\SO(n)/\SO(n-2)$ in \cite{Abiev25b}, where the isotropy representations decompose into a one-dimensional trivial summand and two equivalent summands.
% He showed that every $\SO(n)$-invariant metric evolves to one with positive Ricci curvature and that the set of $\SO(n)$-invariant metrics with positive Ricci curvature is forward-invariant. 
% In \cite{DGAZpreprint}, DeVito, Gonz\'alez-\'Alvaro, and Zarei showed that Ricci flow preserves positive sectional curvature on homogeneous spheres and complex projective spaces.
% For more results on homogeneous Ricci flow, see the work of Abiev \cite{Abiev25a}, Abiev et.~al.~\cite{AANS14}, B\"ohm \cite{Bohm15}, B\"ohm and Lafuente \cite{BohmLafuente18}, Buzano \cite{Buzano14}, Cavenaghi et.~al.~\cite{CGM}, Gonz\'alez-\'Alvaro, and Zarei \cite{GAZ24,GAZ25}, Lauret \cite{Lauret13}, Sbiti \cite{Sbiti20}, and Statha \cite{Statha22}.

In this paper, we investigate the preservation of positive Ricci curvature for $\gG$-invariant Riemannian metrics under normalized Ricci flow on simply connected homogeneous spaces $\gG/\gH$ where $\gG$ is compact and semisimple, $\gH$ is closed and connected, and the isotropy representation splits into \textit{exactly two equivalent} irreducible summands. 
In the case that $\gG$ is simple, it follows from the classification by Dickinson and Kerr in \cite{DickinsonKerr08} that $\gG/\gH$ is $\Spin(8)/\gG_2$, which is diffeomorphic to $S^7 \times S^7$.
Our main result is the following:

\begin{result}\label{mainthm:Ricpos}
        Every $\Spin(8)$-invariant metric on $\Spin(8)/\gG_2$ evolves to a metric with positive Ricci curvature under normalized Ricci flow.
        Furthermore, the family of $\Spin(8)$-invariant metrics with positive Ricci curvature is forward-invariant under the flow.
\end{result}

Alternatively, in the case that $\gG/\gH$ has exactly two irreducible isotropy summands and $\gG$ is \textit{not} simple, Pulemotov and Ziller observed in \cite[Section 5.2]{PulemotovZiller21preprint} that $\gG/\gH$ must be a Ledger-Obata space $(\gH\times \gH\times \gH)/\Delta \gH$, where $\gH$ is simple and $\Delta \gH$ denotes the diagonal embedding of $\gH$ into $\gH^3$. 
They also show that the scalar curvature functional of $\gH^3$-invariant metrics on all Ledger-Obata spaces are proportional to that of $\Spin(8)/\gG_2$.
Since normalized Ricci flow is the $L^2$-gradient flow of the scalar curvature, it follows that the ODE systems of normalized Ricci flow for these spaces all coincide. 
% One can verify this fact by directly computing the Ricci curvatures of $\Spin(8)/\gG_2$ and the Ledger-Obata spaces. 
Consequently, Theorem \ref*{mainthm:Ricpos} immediately yields the following:

\begin{corollaryresult}\label{cor:LedgerObata}
    Let $\gG/\gH$ be a simply connected homogeneous space for which $\gG$ is compact and semisimple, $\gH$ is closed and connected, and the isotropy representation splits into exactly two equivalent irreducible summands.
    Then every $\gG$-invariant metric on $\gG/\gH$ evolves to a metric with positive Ricci curvature under normalized Ricci flow.
    Furthermore, the family of $\gG$-invariant metrics with positive Ricci curvature is forward-invariant under the flow.
\end{corollaryresult}

We note that our result has a conclusion similar to Abiev's main result on the Stiefel manifolds $\SO(n)/\SO(n-2)$ in \cite{Abiev25b}.
While Abiev exploits extra symmetries to diagonalize the $\SO(n)$-invariant metrics on $\SO(n)/\SO(n-2)$ (see \cite[Section 4]{Kerr98}), the family of $\gG$-invariant metrics we consider are not simultaneously diagonalizable.
Consequently, we must analyze the full family of $\gG$-invariant metrics rather than a diagonal subfamily.

Our proof relies on the fact that the family of $\gG$-invariant metrics with fixed volume is two-dimensional in our case.
Thus, the dynamics of the ODE system associated with normalized Ricci flow can be completely visualized.
The spaces we consider here are among the few remaining for which the family of $\gG$-invariant metrics is so low-dimensional while the behavior of normalized Ricci flow has not been studied.
    
\subsection{Organization of paper}

In Section \ref{sec:invariantmetrics}, we give an overview of the isotropy representation, family of $\Spin(8)$-invariant metrics, Ricci curvature, and scalar curvature of $\Spin(8)/\gG_2$.
We then write the normalized Ricci flow equation and highlight the fixed points and several important invariant curves for the system.
In Section \ref{sec:mainproof}, we adopt a more convenient coordinate system for analyzing the system and prove Theorem \ref{mainthm:Ricpos}.
We point out that in several arguments, we rely on the IsEmpty function from Maple’s RegularChains package and SemiAlgebraicSetTools subpackage to prove that certain sets of polynomial inequalities imply others.
The Maple code we used to justify these claims is available on the GitHub repository \cite{GitHub}.
For an overview of this subpackage, see \cite{MapleSite} and the references therein.

\subsection{Acknowledgments}

We would like to thank Jason DeVito, David Gonz\'alez-\'Alvaro, Lee Kennard, Megan Kerr, Yuri\u{i} Nikonorov, Artem Pulemotov, William Wylie, Masoumeh Zarei, and Wolfgang Ziller for helpful discussions. The first author was partially supported by NSF grant DMS-2342135. The second and third authors were supported by Semmes Distinguished Scholars in Science scholarships, and the third author was supported by an AMS-Simons Research Enhancement Grants for PUI Faculty.

    
    % There are two distinct Einstein metrics on $\Spin(8)/\gG_2 \cong S^7 \times S^7$: the product metric, and the metric induced by the Killing form on $\so(8)$.
    % We determine the following about the asymptotics of the system:
    
    % \begin{theorem}\label{mainthm:asymp}
    %     The metric induced by the Killing is an unstable source fixed point for the system.
    %     All other metrics are asymptotically attracted under the flow to one of the following:
    %     \begin{itemize}
    %         \item One of the two fixed points that correspond to the product metric, % \todo{We haven't determined if the one off the ``invariant hyperbola'' serves as an asymptote}
    %         \item The collapsed metric at infinity ($a\to 0^+, b\to \infty$), \todo{This one will be isometric to the orbit space of the appropriate $\Spin(7)$ (or $\SO(7)$?) action on $\Spin(8)/\gG_2$.}
    %         \item The other metric at infinity ($a = 3b$, $a,b \to \infty$).
    %     \end{itemize}
    % \end{theorem}

    % \todo[inline]{Are the metrics on the line $a = 3b$ special in some sense? The Einstein metric this line emanates from is a product metric. Are they all product metrics, maybe with round metrics on the factors with different radii? See Figure \ref{fig:phaseportait}.}

    % \todo[inline]{Are there integral curves that approach the Einstein metric which isn't on the curve $b=1/2a$? That Einstein metric is a product metric, according to Kerr. Maybe there is a curve of product metrics that approach it. Maybe something else entirely... There could even be a integral curve coming from the neighboring Einstein metric.}

\section{Ricci flow equation on the family of invariant metrics}
\label{sec:invariantmetrics}

\subsection{Invariant metrics}
    % As described by Kerr in \cite{Kerr98}, the simple Lie group $\Spin(8)$ acts transitively on $S^7 \times S^7$ such that the isotropy group of the element $(1,1)$ is isomorphic to $\gG_2$.
    By the classification carried out by Dickinson and Kerr in \cite{DickinsonKerr08}, $\Spin(8)/\gG_2$ is the only simply connected homogeneous space $\gG/\gH$ such that $\gG$ is simple and compact, $\gH$ is a connected and closed subgroup, and the isotropy representation splits into exactly two equivalent subrepresentations. 

    To describe the isotropy representation of this space, we will follow \cite{Kerr98,Kerin11}.
    Letting $\bO$ denote the octonions, or Cayley numbers, there is a canonical matrix group representation for $\Spin(8)$: 
    \[
        \Spin(8) = \{(A,B,C) \in \SO(8) : A(x) B(y) = C(xy) \; \text{ for all } x, y \in \bO \}.
    \]
    $\Spin(8)$ acts transitively on the product of spheres $S^7 \times S^7 \subset \bO \times \bO$ via the map $(A,B,C) : (x,y) \mapsto (Ax,By)$, and the isotropy group of the element $(1,1)$ is 
    \[
        \gG_2 = \{ (A,B,C) \in \Spin(8) : A = B = C \}.
    \]
    Thus $\Spin(8)/\gG_2$ is diffeomorphic to $S^7 \times S^7$.

    Now $\Spin(8)$ is the universal cover of $\SO(8)$, with a double covering homomorphism $\Spin(8)\to\SO(8)$ given by $(A,B,C) \mapsto C$.
    Thus, the Lie algebra of $\Spin(8)$ can be identified with $\so(8) = \{A \in M_8(\bR) : A^t = -A\}$.
    For an appropriately ordered basis of $\bO$, we have the inclusions of Lie groups
    \[
        \gG_2 \subset 
        \begin{pmatrix}
            1 & 0 \\
            0 & \SO(7)
        \end{pmatrix}
        \subset \SO(8),
    \]
    which induce inclusions of Lie algebras $\fg_2 \subset \diag(0,\so(7)) \subset \so(8)$.
    We will identify $\fg_2$ and $\so(7)$ with their inclusions into $\so(8)$.
    Hence, $\fg_2$ consists of matrices of the form
    \[
        \left( 
            \begin{array}{c|ccc}
              0 &  & 0 &  \\ \hline
               &  &  &  \\
              0 &  & U &  \\
               &  &  &  
            \end{array} 
        \right),
    \]
    where $U$ is a skew-symmetric matrix given by
    \[
        \begin{pmatrix}
            0                & u_1+u_2 & u_7+u_8 & u_3+u_4 & u_9+u_{10} & u_5+u_6           & u_{11}+u_{12} \\
            *       & 0       & u_{13}  & -u_{11} & u_5        & -u_9              & u_3           \\
            *      & * & 0       & u_6     & u_{12}     & -u_4              & -u_{10}       \\
            *      & *  & *    & 0       & u_{14}     & u_7               & -u_1          \\
            *    & *    & * & * & 0          & u_2               & u_8           \\
            *      & *    & *    & *    & *       & 0                 & u_{13}+u_{14} \\
            * & *    & *  & *     & *       & *  & 0             \\
        \end{pmatrix},
    \]
    where $(u_1, \dots , u_{14}) \in \bR^{14}$.
    Now any bi-invariant metric on $\SO(8)$ is given as a negative real scalar multiple of the Killing form on $\so(8)$.
    We will choose $Q$ to be the bi-invariant metric on $\SO(8)$ given by
    \[
        Q(X,Y) = - \frac{1}{2} \trace(XY) \text{ for all } X,Y \in \so(8),
    \]
    % $Q(X,Y) = - \frac{1}{2} \trace(XY)$
    We then have the $Q$-orthogonal decomposition
    \[
        \so(7) = \fm \oplus \fg_2,
    \]
    where $\fm$ consists of matrices parametrized by $\vec{v} = (v_1,\dots,v_7) \in \bR^7$ of the form
    \begin{equation}
    \label{eq:M}
        M(\vec{v}) = 
        \begin{pmatrix}
            0   & 0       & 0     & 0     & 0     & 0     & 0     & 0     \\
            0   & 0       & v_1   & v_2   & v_3   & v_4   & v_5   & v_6   \\
            0   & -v_1    & 0     & v_7   & v_6   & -v_5  & v_4   & -v_3  \\
            0   & -v_2    & -v_7  & 0     & -v_5  & -v_6  & v_3   & v_4   \\
            0   & -v_3    & -v_6  & v_5   & 0     & v_7   & -v_2  & v_1   \\
            0   & -v_4    & v_5   & v_6   & -v_7  & 0     & -v_1  & -v_2  \\
            0   & -v_5    & -v_4  & v_3   & v_2   & v_1   & 0     & -v_7  \\
            0   & -v_6    & v_3   & -v_4  & -v_1  & v_2   & v_7   & 0     \\
        \end{pmatrix}.
    \end{equation}
    Similarly, we have a $Q$-orthogonal decomposition 
    \[
        \so(8) = \fp \oplus \so(7),
    \]
    where $\fp$ is parametrized by $\vec{w} \in \bR^7$ via
    \begin{equation}
    \label{eq:P}
        P(\vec{w}) = 
        \left( 
            \begin{array}{c|ccc}
              0 &  & -\vec{w}^T &  \\ \hline
               &  &  &  \\
              \vec{w} &  & 0 &  \\
               &  &  &  
            \end{array} 
        \right)
        % \begin{pmatrix}
        %     0   & -w_1  & -w_2  & -w_3  & -w_4  & -w_5  & -w_6  & -w_7  \\
        %     w_1 & 0     & 0     & 0     & 0     & 0     & 0     & 0     \\
        %     w_2 & 0     & 0     & 0     & 0     & 0     & 0     & 0     \\
        %     w_3 & 0     & 0     & 0     & 0     & 0     & 0     & 0     \\
        %     w_4 & 0     & 0     & 0     & 0     & 0     & 0     & 0     \\
        %     w_5 & 0     & 0     & 0     & 0     & 0     & 0     & 0     \\
        %     w_6 & 0     & 0     & 0     & 0     & 0     & 0     & 0     \\
        %     w_7 & 0     & 0     & 0     & 0     & 0     & 0     & 0     \\
        % \end{pmatrix}.
    \end{equation}
    The isotropy representation of $\gG_2$ on $\fg_2^\perp \subset \so(8)$ splits into two equivalent subrepresentations that respect the splitting $\fg_2^\perp = \fm \oplus \fp$.
    Following \cite[Equation 21]{Kerin11}, a $\gG_2$-equivariant linear isomorphism $\psi:\fm \to \fp$ is given by
    \[
        \psi(M(v_1, \dots ,v_7)) = P(v_7, -v_2, v_1, -v_4, v_3, v_6, -v_5).
    \]
    Letting $e_1,\dots,e_7$ denote the standard basis for $\bR^7$, we fix a $Q$-orthonormal basis $\{X_1,\dots,X_{14}\}$ for $\fm \oplus \fp$ given by 
    % {\color{blue}
        \begin{equation}\label{eq:Xi}
            X_i = \frac{1}{\sqrt{3}}M(e_i)\text{ and }X_{i+7} = \psi(M(e_i))\text{ for }1 \leq i \leq 7.
        \end{equation}
    % }
    
    The complexifications of the $\gG_2$-representations on $\fm$ and $\fp$ are irreducible.
    Thus, it follows from Schur's lemma that with respect to the basis $\{X_i,\dots,X_{14}\}$, all $\gG_2$-invariant inner products $g$ on $\fm\oplus\fp$ are of the form 
    \[
        g = \begin{pmatrix}
            a I_7 & c I_7 \\
            c I_7 & b I_7
        \end{pmatrix},
    \]
    where $a,b,c\in\bR$, $a>0$, $b>0$, $ab-c^2 > 0$, and $I_7$ is the $7 \times 7$ identity matrix.
    The family of $\Spin(8)$-invariant metrics on $\Spin(8)/\gG_2$ is in one-to-one correspondence with the set of $\mathrm{Ad}(\gG_2)$-invariant inner products on $\fm\oplus\fp$, so we will refer to an element of either set as simply $g$.
    
    % With respect to this metric, we get orthogonal decompositions 
    % \[
    %     \so(8) = \fp \oplus \diag(1,\so(7)) \qquad \text{and} \qquad .
    % \]
        
    % There are three subgroups intermediate for the pair $\gG_2 \subset \Spin(8)$ that are isomorphic to $\Spin(7)$, and these three are permuted by the triality automorphism of $\Spin(8)$.
    % For more information, see \cite{Kerr98}.

    
    % Because $\Spin(8)$ is the universal cover of $\SO(8)$, its Lie algebra is isomorphic to $\so(8) = \{A \in M_n(\bR) : A^t = -A\}$.
    
    
    
    
    % \todo[inline]{In the exposition above, we will define a basis $\{X_i\}_{i=1}^{14}$ of $\fm \oplus \fp$ that we will use as our reference frame for our space of metrics and system of ODEs.}
    
    % All invariant metrics on $\Spin(8)/\gG_2$ can be written with respect to the basis $\{X_i\}_{i=1}^{14}$ in the form
    % \[
    %     g = \begin{pmatrix}
    %         a I_7 & c I_7 \\
    %         c I_7 & b I_7
    %     \end{pmatrix},
    % \]
    % where $a,b,c$ are real numbers, $a>0$, $b>0$, $ab-c^2 > 0$, and $I_7$ is the $7 \times 7$ identity matrix.

    % \todo[inline]{Will expand upon this later.}

\subsection{Curvature}
    Given an invariant metric $g$ on $\Spin(8)/\gG_2$, the Ricci tensor can be identified with a $G_2$-equivariant endomorphism on $\fm \oplus \fp$. 
    Thus, it also follows from Schur's lemma that the Ricci $(0,2)$-tensor for $g$ can be represented in the basis $\{X_1, \dots, X_{14}\}$ by
    \[
        \of{\Ric_{i,j}} = \begin{pmatrix}
            \alpha I_7 & \gamma I_7 \\
            \gamma I_7 & \beta I_7
        \end{pmatrix}, 
    \]
    for some real numbers $\alpha, \beta, \gamma$.
    By adapting \cite[Corollary 7.33]{Besse87}, or using \cite[Lemma 3.6]{Puttmann99}, it is straightforward to compute that 
    % {\color{blue}
        \begin{align*}
            \alpha &= \frac{9}{2} + \frac{(a^2 - 3 c^2)^2}{2(a b - c^2)^2},\\
            \beta &= \frac{-2 a^3 b + 3 a^2 (4 b^2 + c^2) - 24 a b c^2 + 9 b^2 c^2 + 6 c^4}{2(a b - c^2)^2}, \\
            \gamma &= \frac{a c (a - 3 b)^2}{2(a b - c^2)^2}.
        \end{align*}
        % Performing the Gram-Schmidt process on the vectors $\{X_i\}_{i=1}^{14}$ defined in \eqref{eq:Xi}, a $g$-orthonormal basis $\{E_i\}_{i=1}^{14}$ of $\fm\oplus\fp$ is given by 
        % \[
        %     E_i = \tfrac{1}{\sqrt{a}} X_i \text{ and } E_{i+7} = -\sqrt{2b-\tfrac{1}{a}} \; X_i + \sqrt{2a} \; X_{i+7} \text{ for } 1 \leq i \leq 7.
        % \]
        % With respect to this basis, the Ricci $(1,1)$-tensor is given by 
        Then in the basis $\{X_1, \dots, X_{14}\}$, the Ricci $(1,1)$-tensor can be represented by
        \[
            \of{\Ric^i_j} = \of{g^{i,k} \Ric_{k,j}} = 
            \begin{pmatrix}
                \delta I_7 & \epsilon I_7 \\
                \zeta I_7 & \eta I_7
            \end{pmatrix}
        \]
        where
        \begin{equation}\label{eq:Ric11abc}
        \begin{aligned}
            \delta &= \frac{a^3 + 9 a b^2 - 18 b c^2}{2 (a b - c^2)^2}, \\
            \epsilon &= \frac{3c (a^2 - 6 a b + 3 b^2 + 2 c^2)}{2 (a b - c^2)^2}, \\
            \zeta &= -\frac{3c (a^2 - 3 c^2)}{(a b - c^2)^2},\\
            \eta &= -\frac{a (a^2 - 6 a b + 3 c^2)}{(a b - c^2)^2}.
        \end{aligned}
        \end{equation}
        
        % \begin{align}
        %     p &= \frac{9+(2a^2-6ab+3)^2}{2a}, \label{eqn p}\\
        %     q &= -\frac{(16a^4-96a^3b+9(4ab-2)^2)}{4a}, \label{eqn q}\\
        %     r &= -\frac{3\sqrt{2ab-1}(2a^2-6ab+3)}{a}.\label{eqn r}
        % \end{align}
        This operator has two real eigenvalues given by the roots of the polynomial $(\delta-x)(\eta-x)-\epsilon\zeta$.
        To show that these roots are positive, it is enough to check that their sum $S=\delta+\eta$ and product $P=\delta\eta-\epsilon\zeta$ are both positive.
        Thus, we have
        \begin{align*}
            S &= -\frac{a^3 - 12 a^2 b - 9 a b^2 + 6 a c^2 + 18 b c^2}{2 (a b - c^2)^2}\\
            P &= \frac{-a^6 + 6 a^5 b - 9 a^4 b^2 + 6 a^4 c^2 + 54 a^3 b^3 - 36 a^3 b c^2 }{2 (a b - c^2)^4} \\
            &\phantom{==} + \frac{- 108 a^2 b^2 c^2 - 9 a^2 c^4 + 216 a b c^4 - 81 b^2 c^4 - 54 c^6}{2 (a b - c^2)^4}.
        \end{align*}
        % By Vieta's formulas, 
        % \begin{align}\label{eqn P}
        %     S &= p + s,\nonumber \\
        %     P &= p q - r^2.
        % \end{align}
        
        Finally, it follows that the scalar curvature is
        \[
            \Scal = -\frac{7(a^3 - 12 a^2 b - 9 a b^2 + 6 a c^2 + 18 b c^2)}{2 (a b - c^2)^2}.
        \]
    % }

    

\subsection{Normalized Ricci flow}
    The normalized Ricci flow equation for an $n$-dimensional manifold $M^n$ with initial metric $g_0$ is as follows:
    \[
        \frac{\partial}{\partial t} g(t) = -2 \Ric_{g(t)} + \frac{2 \Scal}{n} g(t), \qquad g(0)=g_0.
    \]
    % {\color{blue}
        Applying our computations above, we have the following system of equations for the normalized homogeneous Ricci flow on $\Spin(8)/\gG_2$:
        \begin{align*}
            &
            a'(t)=-\frac{3(a^4 - 4 a^3 b + 3 a^2 b^2 - 2 a^2 c^2 - 6 a b c^2 + 12 c^4)}{2 (a b - c^2)^2},\\
            &
            b'(t)=\frac{9 a b^3 - 12 b^2 ( a^2 + 3 c^2)  + 3 b (a^3 + 14 a c^2) - 6 c^2 (a^2 + 2 c^2)}{2 (a b - c^2)^2},\\
            &
            c'(t)=-\frac{3c (a^3 - 8 a^2 b + 3 a b^2 + 2 a c^2 + 6 b c^2)}{2 (a b - c^2)^2},\\
            & 
            a>0, \quad b>0, \quad ab-c^2>0.
        \end{align*}
    % }
    
    We will normalize the volume of $\Spin(8)/\gG_2$ such that 
    \begin{equation}\label{eq:volume_normalization}
        ab-c^2 = 1/2.
    \end{equation}
    Then normalized Ricci flow equation is given in the variables $a,b$ as follows:
    \begin{equation}\label{eq:sys_ab}
    \begin{aligned}
        &
        a'(t) = -18 - 6 a^4 + 36 a^3 b + (-54 b^2 - 6) a^2 + 54 a b,  \\
        &
        b'(t) = -6 a^3 b + (36 b^2 + 6) a^2 + (-54 b^3 - 18 b) a + 36 b^2 - 6,  \\
        &
        a>0, \quad b>0, \quad 2ab \geq 1.
    \end{aligned}
    \end{equation}
    We will call $\{(a,b)\in\bR^2:a>0,b>0,2ab\geq 1\}$ the set of admissible metrics.

    Fixed points for the flow occur at 
    \begin{equation}\label{eq:sing_ab}
        (a,b) = \of{ \tfrac{\sqrt{2}}{2} , \tfrac{\sqrt{2}}{2} }, \of{ \tfrac{\sqrt{6}}{4} , \tfrac{5\sqrt{6}}{12} }, \text{ and } \of{ \tfrac{\sqrt{6}}{2} , \tfrac{\sqrt{6}}{6} }.
    \end{equation}
    Each of these fixed points corresponds to a homogeneous Einstein metric on $\Spin(8)/\gG_2$ found by Kerr in \cite{Kerr98}.
    The metric induced by the Killing form on $\so(8)$ corresponds to $( \frac{\sqrt{2}}{2} , \frac{\sqrt{2}}{2} )$, while a product metric on $S^7 \times S^7 \cong \Spin(8)/\gG_2$ corresponds to $( \frac{\sqrt{6}}{4} , \frac{5\sqrt{6}}{12} )$.
    Finally, the metric corresponding to $(\frac{\sqrt{6}}{2} , \frac{\sqrt{6}}{6} )$ can be obtained from this product metric by conjugating by a rotation matrix:
    \[
        \begin{pmatrix}
            \frac{1}{2}         & -\frac{\sqrt{3}}{2} \\
            \frac{\sqrt{3}}{2}  & \frac{1}{2}
        \end{pmatrix}
        \begin{pmatrix}
            \frac{\sqrt{6}}{2}  & 0 \\
            0                   & \frac{\sqrt{6}}{6}
        \end{pmatrix}
        \begin{pmatrix}
            \frac{1}{2}         & \frac{\sqrt{3}}{2} \\
            -\frac{\sqrt{3}}{2} & \frac{1}{2}
        \end{pmatrix}
        =
        \begin{pmatrix}
            \frac{\sqrt{6}}{4}      & \frac{\sqrt{18}}{12} \\
            \frac{\sqrt{18}}{12}    & \frac{5\sqrt{6}}{12}
        \end{pmatrix}        
    \]
    As Kerr writes in \cite{Kerr98}, this rotation is the action of the triality automorphism of $\so(8)$, and the metric corresponding to $( \frac{\sqrt{6}}{2} , \frac{\sqrt{6}}{6})$ is isometric to a product metric.
    

    % \todo[inline]{Indicate how this is an easy check. Talk about attracting and repelling manifolds for fixed points. Say how we found one from another using transformation in Kerr '98.}
    
    % \begin{proposition}\label{prop:invariant_curves_ab}
        It is straightforward to check that the following are invariant curves for the flow of the System \eqref{eq:sys_ab}:
        \begin{enumerate}
            \item $2ab = 1$,                        % Boundary curve
            \label{curve1ab}
            
            \item $a = 3b$,                         % Other curve going through (sqrt(6)/2, sqrt(6)/6)
            \label{curve2ab}
            
            \item $3( a^2 + b^2 ) + 2 = 10 a b$.    % Attracting curve for (sqrt(6)/4, 5sqrt(6)/12)
            % $\{ a = \frac{\tau}{4} + \frac{3}{8\tau}, b = \frac{3 \tau}{4} + \frac{1}{8 \tau}\}$, where $\tau > 0$.
            \label{curve3ab}
            
            \item $2a^2 + 3 = 6ab$.                 % Repelling curve for (sqrt(6)/4, 5sqrt(6)/12)
            \label{curve4ab}
        \end{enumerate}
    % \end{proposition}
    (\ref*{curve1ab}) is the boundary curve for the set of admissible metrics, which corresponds to where $c=0$.
    We note that (\ref*{curve4ab}) can be obtained from (\ref*{curve2ab}) by applying the same transformation that sends $\of{ \frac{\sqrt{6}}{2} , \frac{\sqrt{6}}{6} }$ to $\of{ \frac{\sqrt{6}}{4} , \frac{5\sqrt{6}}{12} }$.
    
    Figure \ref{fig:phaseportait} illustrates a phase portrait for System \eqref{eq:sys_ab} on the left.
    The gray region consists of inadmissible values of $(a,b)$ for which $2ab<1$, the red points are fixed points for the system which correspond to Einstein metrics, the black dotted curves are the invariant curves we have listed above, and the green region is the set of metrics for which $\Ric>0$.
    
    

% \subsection{Invariant curves of the system}

    
    \begin{figure}[H]
        \centering
        \includegraphics[width=0.49\textwidth]{DEplot.eps}\includegraphics[width=0.49\textwidth]{DEplotXY.eps}
        \caption{Phase portraits for Systems \eqref{eq:sys_ab} and \eqref{eq:sys_xy}.}
        \label{fig:phaseportait}
    \end{figure}
    

\section{Positive Ricci curvature in the long term}
\label{sec:mainproof}

    % \todo[inline]{Lawrence's goals for this section: Explain the road map of what we show in this section ($P>0$ is forward invar. and invariant curves eventually have $P>0$, contained in a compact works, $|x|$ bounded works, then we finish remaining case.) with a picture, and adapt proof of the first few lemmas to work in $xy$ coordinates.}

    

    In this section, we will prove Theorem \ref{mainthm:Ricpos}.
    Namely, we will show that all solutions to the System \eqref{eq:sys_ab} eventually develop and maintain positive Ricci curvature.

    Because the line $a=3b$ is an invariant curve of the System \eqref{eq:sys_ab}, we find it convenient to implement the change of coordinates $x= a - 3b$ and $y=b$.  In these new coordinates, System \eqref{eq:sys_ab} becomes:
    \begin{equation}\label{eq:sys_xy}
    \begin{aligned} 
        &
        x'(t) = -6x\left(x^3 + 3x^2 y + 4x + 6y\right), \\
        &
        y'(t) = -6x^3 y - 18x^2 y^2 + 6x^2 + 18x y + 36y^2 - 6. \\
        &
        2xy+6y^2\geq 1, \quad y>0  .
    \end{aligned}
    \end{equation}
    Figure \ref{fig:phaseportait} illustrates a phase portrait for System \eqref{eq:sys_xy} on the right. 
    The fixed points from \eqref{eq:sing_ab} are given by
    \begin{equation}\label{eq:sing_xy}
        (x,y) = \of{-\sqrt{2},\tfrac{\sqrt{2}}{2}}, \of{-\sqrt{6},\tfrac{5\sqrt{6}}{12}}, \text{ and } \of{0,\tfrac{\sqrt{6}}{6}}.
    \end{equation}
    The invariant curves from \eqref{curve1ab}--\eqref{curve4ab} above are given in the $(x,y)$-coordinates by
    \begin{enumerate}
        \item \label{curve1xy} $2xy + 6y^2 = 1$ (with $y>0$),
        \item \label{curve2xy} $x=0$,
        \item \label{curve3xy} $3x^2 + 8xy + 2 = 0$ (with $x<0$), and %$3x^2 + 8xy + 2 = 0$,
        \item \label{curve4xy} $2x^2 + 6xy + 3 = 0$ (with $x<0$). %$2x^2 + 6xy + 3 = 0$.
    \end{enumerate}
    The entries of the Ricci $(1,1)$-tensor from \eqref{eq:Ric11abc} become
    \begin{align*}
        \delta &= 2 x^3 + 18 x^2 y + 36 x y^2 + 18 y, \\
        \epsilon &= 3 (x^2 + 2 x y - 1) \sqrt{4 x y + 12 y^2 - 2} , \\
        \zeta &= -3 (2x^2 + 6 x y + 3) \sqrt{4 x y + 12 y^2 - 2} ,\\
        \eta &= -2 (2x^2 + 6 x y - 3) (x + 3 y).
    \end{align*}
    Thus, the sum $S$ and product $P$ of the eigenvalues of the Ricci tensor (counted without multiplicity) are 
    \begin{equation}\label{eq:Pxy}
    \begin{aligned}
        S &= -2 x^3 - 6 x^2 y + 6 x + 36 y,  \\
        P &= 54 - 8 x^6 - 48 x^5 y + (-72 y^2 - 24) x^4 - 72 x^3 y - 18 x^2.
    \end{aligned}
    \end{equation}
    %     \begin{align}
    % p &= \frac{9+(2x^2+6xy+3)^2}{2(x+3y)}, \\ 
    % q &= -\frac{16(x+3y)^4 - 96(x+3y)^3 y + 9(4xy+12y^2-2)^2}{4(x+3y)}, \\[6pt]
    % r &= -\frac{3\sqrt{2xy+6y^2-1}\,(2x^2+6xy+3)}{x+3y},\\[6pt]
    % P &= -8x^6 - 48x^5y - 72x^4y^2 - 24x^4 - 72x^3y - 18x^2 + 54.\label{eq:Pxy}\\
    % \end{align}
    
    
    Throughout this section, $(x(t),y(t))$ will refer to a solution to the System \eqref{eq:sys_xy}, and $g(t) = g(x(t),y(t))$ will denote the associated metrics.
    When discussing aspects of these objects that are not time dependent, we will refer to them simply as $(x,y)$ and $g$, respectively.
    
    
    We first show that it is enough to check for the positivity of $P$:

    % \todo[inline]{This next proof is from before we realized there is an error in how we were writing $P$. The formula above for $P$ in terms of $x$ and $y$ is correct, but it is not $pq-r^2$ as in the proof. Also, the terms in $S$ used in the proof may be old and incorrect. Need to come back and fix this.}
    
    \begin{lemma}\label{lem:ISTS_P>0}
        A metric $g$ has $\Ric > 0$ if and only if $P > 0$.
    \end{lemma}

    \begin{proof}
        %It is clear that $\Ric> 0$ implies $P>0$.
        %Conversely, if $a,b,P>0$, then it follows that $S>0$.

        %Eric's attempt at a proof 10/1/25 ------------------

        Because $P$ is the product of eigenvalues of the  Ricci tensor, it is clear that $\Ric> 0$ implies $P>0$. 
        Conversely, using Maple's RegularChains package and SemiAlgebraicSetTools subpackage, the IsEmpty function shows that $\{(x,y) : y>0, \ 2 x y + 6 y^2 \geq 1, \ P > 0, \ S \leq 0\}$ is empty.
        Thus if $P>0$, then $S>0$, and because the Ricci tensor has only two eigenvalues (counted without multiplicity), it follows that $\Ric>0$.
        % Conversely, let $P > 0$.  
        % Then $pq = r^2 + P > 0$ and so $p$ and $q$ have the same sign.  
        % By assumption, $y>0$ and $2xy+6y^2\geq 1$, which implies $x+3y>0$.
        % Thus, it follows from the expression for $p$ above that $p > 0$, and so necessarily $q>0$ and $S = p + q > 0$ as well.  
        % Therefore, as observed above, it follows that $\mathrm{Ric}>0$.
        % 
        % $\Ric(v,v) = \Ric(\alpha^i v_i, \alpha^j v_j) = \alpha^i \alpha^j \Ric(v_i,v_j) = \alpha^i \alpha^j \Ric(v_i,v_j) \delta_{i,j} = (\alpha^i)^2 \Ric(v_i,v_i) > 0$
        % 
        %----------------------------------------------------
        % 
        % \todo[inline]{Arseny and Lawrence checked this using Maple. We can reference that and provide code online, or we can try to prove this by hand...}
    \end{proof}

    \begin{remark}\label{rmk 1}
        Setting $P$ from \eqref{eq:Pxy} equal to zero and solving for $y$ yields:
        \[
        y_1(x) = -\frac{x}{3} - \frac{1}{2x} - \frac{\sqrt{3}}{2x^2},
        \qquad
        y_2(x) = -\frac{x}{3} - \frac{1}{2x} + \frac{\sqrt{3}}{2x^2}.
        \]
        Notice $y_1<y_2$ for all $x\neq0$.
        Since the coefficient of $y^2$ is negative in \eqref{eq:Pxy}, we have that $P$ is positive only when $y_1<y<y_2$. 
        Furthermore, if $y>0$ and $2xy+6y^2\geq 1$, then using Maple's IsEmpty function, one can show $y_1 < y$.
        Therefore, we have $P>0$ if and only if $y<y_2$.
    \end{remark}

    Next, we show that when composed with a solution of system \eqref{eq:sys_xy}, $P$ is a strictly increasing function of $t$ almost everywhere.
    First, denoting $P|_{(x(t),y(t))}$ as simply $P(t)$, we remark that
    \begin{align*}
        P'(t) &= 72x^2\left(2x^2 + 6xy + 3\right)^2 \left(x^3 + 3x^2y + 2x + 2y \right).
    \end{align*}
    
    \begin{lemma}
        \label{lem:Pinc}
        For any $(x,y)$ lying in the region where $2xy+6y^2>1$ and $y>0$, we have $P'(t)>0$ for all $t$ except on the two invariant curves $x=0$ and $2x^2+6xy+3=0$ (with $x<0$), on which $P'(t)$ is zero. 
        Additionally, these two curves are contained in the region where $P>0$.
        Consequently, the region where $P>0$ is forward-invariant.
    \end{lemma}
    
    \begin{proof}
        First, it is again verifiable using Maple's IsEmpty function that $P'(t)\geq 0$ when $2xy+6y^2 > 1$ and $y>0$. 
        Because $x$ and $2x^2+6xy+3$ are factors of $P'(t)$, the curves $x=0$ and $2x^2+6xy+3=0$ satisfy $P'(t)=0$.
        The zeros of the remaining factor $x^3 + 3x^2y + 2x + 2y$ of $P'(t)$ lie outside of the set where $2xy+6y^2 > 1$, meaning the invariant curves $x=0$ and $2x^2+6xy+3=0$ constitute the only points where $P'(t)=0$.
        By uniqueness of solutions, it follows that any solution whose initial condition satisfies $2xy+6y^2 > 1$ and lies outside of these curves must have $P'(t)>0$ for all $t$. 
        The claim that all points on these invariant curves satisfy $P>0$ is also verifiable with Maple.
    \end{proof}  
    
    
    The fixed points of System \eqref{eq:sys_xy} given in \eqref{eq:sing_xy} are all contained in the region where $P>0$, and by Lemma \ref{lem:Pinc}, so are the invariant curves \eqref{curve2xy} and \eqref{curve4xy}.
    The other invariant curves \eqref{curve1xy} and \eqref{curve3xy} with the fixed points deleted consist of components that either are contained in the region where $P>0$ or serve as a component of the stable manifold for one of the saddle point singularities at $(x,y)=(-\sqrt{6},\frac{5\sqrt{6}}{12})$ and $(0,\frac{\sqrt{6}}{6})$.
    In particular, we have verified that the conclusion of Theorem \ref{mainthm:Ricpos} holds for initial conditions lying on any of the curves \eqref{curve1xy}--\eqref{curve4xy}.
    
    Due to this fact, below we only work with initial conditions which lie in the complement to the union of these invariant curves.
    Within the interior of the space of admissible metrics, $M=\{(x,y): y > 0 \text{ and } 2xy+6y^2 > 1\}$, this complement is partitioned into the following connected components:
    \begin{equation}\label{regions}
    \begin{aligned}
        A &= \{(x,y)\in M : 3x^{2}+8xy+2 > 0 \text{ and } 2x^2 + 6xy + 3 < 0 \},  \\
        B &= \{(x,y)\in M : x < -\sqrt{2}, 3x^{2}+8xy+2 > 0, \text{ and } 2x^2 + 6xy + 3 > 0 \},   \\
        C &= \{(x,y)\in M : 3x^{2}+8xy+2 < 0  \text{ and } 2x^2 + 6xy + 3 < 0 \},   \\
        D &= \{(x,y)\in M : 3x^{2}+8xy+2 < 0 \text{ and } 2x^2 + 6xy + 3 > 0 \},   \\
        E &= \{(x,y)\in M : -\sqrt{2} < x < 0 \text{ and } 3x^{2}+8xy+2 > 0 \},   \\
        F &= \{(x,y)\in M : x > 0 \}.   
    \end{aligned}
    \end{equation}
    These components are illustrated in Figure \ref{fig:regions}.
    Because they are bounded by invariant curves, they are forward-invariant for the flow of \eqref{eq:sys_xy}.
    Thus, it suffices to prove Theorem \ref{mainthm:Ricpos} for solutions whose initial conditions lie in each of these regions.

    \begin{figure}[H]
        \centering
        \includegraphics[width=0.5\linewidth]{plot_xy_region_A-F.eps}
        \caption{Regions $A$--$F$ defined in \eqref{regions}.}
        \label{fig:regions}
    \end{figure}

    % Next, we establish Theorem \ref{mainthm:Ricpos} for bounded solutions of system \eqref{eq:sys_ab}:
    
    % \begin{lemma}\label{lem:compact}
    %     If $(x(t),y(t))$ is contained in a compact region for all $t$, then there exists $t_0$ such that $g(t)$ has $\Ric > 0$ for all $t \geq t_0$.
    % \end{lemma}

    % \begin{proof}
    %     Suppose for the sake of contradiction that $(x(t),y(t))$ is contained in a compact set $K$ which is a subset of the closed region  where $P \leq 0$.
    %     % Because $P\leq 0$ is a closed condition, we may assume that $K$ is contained in the region where $P\leq 0$.
    %     We will show that there exists a point in $K$ for which $P'(t) = 0$, contradicting Lemma \ref{lem:Pinc}.
        
    %     Because the solution is contained in a compact region, $(x(t),y(t))$ is defined for all $t$.
    %     Define $U = \sup_t P(t)$.
    %     From Lemma \ref{lem:Pinc}, it follows that $P(t) < U$ for all $t$.
    %     {\color{blue}
    %         Thus $U-P(t)$ can be may arbitrarily small for large enough values of $t$, and similarly, so can
    %         \[
    %             \frac{|P(t+2)-P(t)|}{2} \leq \frac{|P(t+2)-U|}{2}+\frac{|U-P(t)|}{2}.
    %         \]
    %         So by the Mean Value Theorem, we can construct an increasing sequence $t_i$ such that $\frac{dP}{dt}|_{t=t_i} \to 0$.
    %     }
    %     {\color{red}
    %         So for all $\ep > 0$, there exists $t_\ep$ such that for all $t \geq t_\ep$, we have $U-P(t) < \ep$.
    %         So for any $t \geq t_\ep$, we get 
    %         \[
    %             \frac{|P(t+2)-P(t)|}{2} \leq \frac{|P(t+2)-U|}{2}+\frac{|U-P(t)|}{2} < \ep.
    %         \]
    %         Thus by the Mean Value Theorem, there exists $s\in[t,t+2]$ such that $\frac{dP}{dt}|_{t=s} < \ep$.
    %         Consequently, we can construct an increasing sequence $t_i$ such that $\frac{dP}{dt}|_{t=t_i} \to 0$.
    %     }
    %     Because the sequence $(x(t_i),y(t_i))$ is contained in the compact region $K$, it has a subsequence that converges to a point $(x_\infty,y_\infty) \in K$.
    %     Since the region where $P\leq 0$ is closed, we have that $(x_\infty,y_\infty)$ satisfies $P\leq0$, but by the continuity of $\frac{dP}{dt}$, it follows that $\frac{dP}{dt}|_{(x_\infty,y_\infty)} = 0$.
    %     %
    %     % \todo[inline]{Need to show: For every $\epsilon > 0$, there exists some $\sigma \in (0,S)$ such that for all $t \in (S - \sigma,S)$, $P'(t) < \epsilon$.}
    %     %
    %     % \todo[inline]{Need to show: For every $\epsilon > 0$, there exists some $\sigma \in (0,S)$ such that for all $t \in (S - \sigma,S)$, there exists some $\eta \in (0,\sigma)$ such that for all $h \in (0,\eta)$, we have $| (P(t+h) - P(t)) / h | < \epsilon$.}
    %     %
    %     % \todo[inline]{Need to show: For every $\epsilon > 0$, there exists some $T > 0$ such that for all $t \geq T$, there exists some $\eta > 0$ such that for all $h \in (0,\eta)$, we have $| (P(t+h) - P(t)) / h | < \epsilon$.}
    %     %
    %     % \todo[inline]{This is the line separating new (above) from old (below).}
    %     %
    %     % Suppose that there exists an initial $(a_0,b_0)$ s.t. the integral curve starting from it never reaches the $P>0$ region. Then the integral curve cannot be contained in a compact region $K$. Let us denote the compact intersection of $K$ with the closed region where $a\geq0,ab\geq\frac{1}{2},P\leq0$ by $K_0$. $P$ has values 25 and 27 at the three fixed points, so we need not worry about them. By our assumption, the integral curve lies completely in $K_0$. Combined with the fact that our equation is smooth, it follows that the solutions can be extended along the $t$ axis indefinitely. We claim that $$\lim_{t\to\infty}\frac{dP}{dt}=0.$$ Indeed, we assume that, while $P(t)$ strictly increases as $t\to\infty$, $$U=\lim_{t\to\infty}P(t)=\sup_t P(t)\leq 0.$$ $U$ is never obtained, since, if it were obtained at $t=t_0$, then beyond that point $P(t)$ would be strictly greater than $U$, contradicting its definition. Therefore, for any $\varepsilon'>0$ there exists a $T$ s.t. for $U-P(T)<\varepsilon'$, and then, by monotonicity, the same inequality holds for all $t>T$. 
    %     % Therefore for $t\geq T$, $$|P(t+\delta)-P(t)|=|P(t+\delta)-U+U-P(t)|<2\varepsilon'$$
    %     % \[
    %     %     |P(t+\delta)-P(t)|/\delta < 2 \epsilon' / \delta
    %     % \]for any $\delta>0$. 
    %     % Therefore, we can pick $\varepsilon'=\frac{\delta\varepsilon}{2}$, so that $$\Big|\frac{P(t+\delta)-P(t)}{\delta}\Big|<\varepsilon$$ for $t>T'$. Since $\varepsilon$ can be chosen to be arbitrarily small, the statement is proven. Next, we would be able to pick a sequence of points $K_0\ni x_n=(a(n),b(n)),n\in\mathbb{N}$ along the integral curve, where $x_n$ are distinct since $P(t)$ strictly increases along them, and a converging subsequence $\{x_{n_k}\},k\in\mathbb{N}$ s.t. $$\lim_{k\to\infty}x_{n_k}=x\in K_0,$$ where for any $\varepsilon>0$ there exists such a $k(\varepsilon)$ that $$\frac{dP}{dt}(x_{n_k})<\varepsilon$$ for all $k\geq k(\varepsilon)$. By the continuity of $\frac{dP}{dt}\in\mathbb{R}[a,b]$ and the previous inequalities, we have $$\lim_{k\to\infty}\frac{dP}{dt}(x_{n_k})=\frac{dP}{dt}(x)=0.$$ This is impossible since $x\in K_0$, but we already established that a point that lies in the region where $a\geq0, ab\geq\frac{1}{2},P\leq 0$ cannot be a zero of $\frac{dP}{dt}.$ Therefore, asymptotic integral curves cannot be contained in a compact region. 
    % \end{proof}

    \begin{remark}\label{rem:compact}
        Note that by Lemma \ref{lem:Pinc}, System \eqref{eq:sys_xy} has no cyclic solutions, as such a solution would necessarily have $P'(t) \equiv 0$, and this equality only holds on the unbounded curves \eqref{curve2xy} and \eqref{curve4xy}.
        Consequently by the Poincar\'e-Bendixson theorem, because the fixed points of the system consist of a source at $(x,y)=(-\sqrt{2},\frac{\sqrt{2}}{2})$ and saddle points at $(x,y)=(-\sqrt{6},\frac{5\sqrt{6}}{12})$ and $(0,\frac{\sqrt{6}}{6})$, the only solutions which are contained in a compact region are the fixed points and solutions on the curves which serve as the stable manifolds for the saddle points.
        As stated above, these stable manifolds lie on curves \eqref{curve1xy} or \eqref{curve3xy}.
        Therefore, because we only need to work with initial conditions which do not lie on these invariant curves, we may assume all solutions analyzed below are unbounded.
    \end{remark}
    % \begin{remark}

    It can be verified with Maple's IsEmpty function that the regions $B$, $D$, and $E$ are contained in the region where $P>0$.
    Thus, we have the following:
    
    \begin{proposition}
        If $(x(t),y(t))$ is a solution of System \eqref{eq:sys_xy} whose initial condition lies in region $B$, $D$, or $E$ from \eqref{regions}, then the solution has $\Ric>0$ for all $t$.
    \end{proposition} 
    
    
    % It also follows from Lemma \ref{lem:Pinc} that the system has no cyclic solutions, as such a solution would necessarily have $P'(t) \equiv 0$, and this equality only holds on the unbounded curves \ref{curve2xy} and \ref{curve4xy}.
    % Consequently by the Poincar\'e-Bendixson theorem, because the fixed points of the system consist of a source at $(x,y)=(-\frac{\sqrt{2}}{2},\frac{\sqrt{2}}{2})$ and saddle points at $(x,y)=(-\sqrt{2},\frac{\sqrt{2}}{2})$ and $(0,\frac{\sqrt{6}}{6})$, it follows that only solutions which are contained in a compact region are the curves which serve as the stable manifolds for the saddle point fixed points.
    % In particular, all of the invariant curves \eqref{curve1xy}--\eqref{curve4xy} consist of 
    
    % In addition to the two invariant curves $x=0$ and 
    % $2x^2+6xy+3=0$ being contained in the region where $P>0$ by Lemma \ref{lem:Pinc}, for any 
    % initial condition chosen on the invariant curves $2xy+6y^2=1$ or 
    % $3x^2+8xy+2=0$, the corresponding solution to System \eqref{eq:sys_xy} will enter the region where $P>0$ in finite time and will 
    % remain there indefinitely.
    % In particular, all points on these curves are either contained in the region where $P>0$, or they are on a component of the curve that serves as a stable manifold for the fixed points $(x,y)=(-\sqrt{2},\frac{\sqrt{2}}{2})$ or $(0,\frac{\sqrt{6}}{6})$, both of which are contained in the region where $P>0$.
    % In particular, we have verified that the conclusion of Theorem \ref{mainthm:Ricpos} holds for initial conditions lying on any of these curves.
    % Due to this fact, below we only work with initial conditions which 
    % lie in the complement to the union of these four invariant curves.    
    % \end{remark}
    

    % \begin{lemma}
    %     The region where $S > 0$ is forward-invariant.
    %     Furthermore, every integral curve must eventually enter $S$.
    % \end{lemma}

    % \begin{proof}
    %     FORWARD INVARIANCE IS CHECKED WITH MAPLE.

    %     In the region where $S \leq 0$, then $\frac{d}{dt}(a^2 + b^2) = 2aa'+2bb' < 0$, AS VERIFIED BY MAPLE.
    %     Thus, if an integral curve $(a(t),b(t))$ were contained in the region where $S \leq 0$, then it would also be contained in any closed ball centered at the origin that also contains the initial condition $(a(0),b(0))$.
    %     Then by Lemma \ref{lem:compact}, the integral curve must enter the region where $P>0$, which implies $S>0$, a contradiction.
    %     Therefore, every integral must eventually enter $S$.
    % \end{proof}

    % \todo[inline]{Let's use $x = a - 3b$ and $y=b$ as our new coordinates. (So Arseny's $\varepsilon = x$ and $b=y$.}

    

    

    % \todo[inline]{Is $3x^2y + x^3  + 2x + 2y$ an invariant curve for our system? Answer: No, this curve is outside the set of admissable metrics.}

    % \todo[inline]{Do our proofs below imply that these invariant curves enter the $P>0$ region, or should we add a remark here?}

    % A consequence of Lemma \ref{lem:ISTS_P>0} is the following:

    
    % \todo[inline]{This cor shows if you start in one of our three regions that 's green, you stay in the green. Add a ``proof.'' Maybe we need forward invariance of $P>0$ before this. Maybe move this right after Lemma 3.2. Make it a Remark?}
    
    
    % \begin{corollary}
    %     If the initial conditions of a solution $(x(t),y(t))$ to System \eqref{eq:sys_xy} satisfy 
    %     \[
    %         y \leq -\frac{2x^2+3}{6x} \quad \text{and} \quad x<0
    %     \]
    %     then the solution satisfies $P>0$ for all values of $t$.
    % \end{corollary}
        
    
    
    Our next goal is to prove Theorem \ref{mainthm:Ricpos} for the remaining solutions that approach the line $x=0$, i.e. those contained in $C$ and $F$.
    We begin with the following:

    \begin{lemma}\label{lemma D}
        Suppose $(x(t),y(t))$ is a solution of System \eqref{eq:sys_xy} whose initial condition lies in region $C$ or $F$ from \eqref{regions}.
        Then $|x(t)|$ is monotonically decreasing for all $t$ and $x(t)\to 0$ and $t$ increases, and $y(t)$ is monotonically increasing for sufficiently large values of $t$ and $y(t) \to \infty$ as $t$ increases.
    \end{lemma}

    % \todo[inline]{Can we prove here if $|x(t)|$ is bounded (no assumption on $P$) and the initial conditions satisfy \eqref{item:condition1} or \eqref{item:condition2} from the next theorem, then $|x(t)|$ converges monotonically to $0$? This would simplify the proof of the next theorem, and it make it so that this lemma is not a vacuous statement.}

    \begin{proof}
        If the initial condition lies in region $F$ (where $x>0$), then $x(t)>0$ for all $t$, and thus $x'(t) = -6x\left(x^3 + 3x^2 y + 4x + 6y\right)<0$ for all $t$.
        If instead the initial condition lies in region $C$, then it can be verified with Maple's IsEmpty function that $x'(t)>0$ for all $t$.
        Therefore, $|x(t)|$ is monotonically decreasing for all $t$ in both cases.
        
        Next, we show that $y(t)$ eventually increases monotonically.
        By Remark \ref{rem:compact}, we may assume that our solution $(x(t),y(t))$ is unbounded.
        Since $|x|$ is bounded, we have that $|y|$ is unbounded.  
        Because $y>0$ by assumption, the solution must satisfy $y'(t) > 0$ for some $t$.  
        Using Maple's IsEmpty function, we can verify if $y>2$, then $y''(t)>0$.
        Thus the region where $y>2$ and $y'(t)>0$ is forward-invariant, and because $y$ is unbounded, our solution must eventually enter and remain in this region.
        Hence, $y(t)$ increases monotonically for all sufficiently large values of $t$, and because $y$ is unbounded, we have $y\to \infty$ as $t$ increases.
        
        

        % Now we show that $x\to 0$ as $t\to\infty$.
        % The region where $y'(t)>0$ has two connected components which are separated by the line $18y+10x=0$; see Figure \ref{fig:y_prime_pos_sep_by_line}.
        % We observe that, because $|x|$ is bounded and $y$ is positive and unbounded, our solution cannot be contained in the region where $18y+10x<0$ and hence must eventually enter and remain in the region where $y>2$, $y'(t)>0$, and $18y+10x>0$.
        % Furthermore, it can be seen directly in \eqref{eq:sys_xy} that $x'(t) < 0$ if $x > 0$, and it can be verified with Maple that $x'(t) > 0$ if $x < 0$, $y>2$, $y'(t)>0$, and $18y+10x>0$.  
        % Since $x=0$ is an invariant curve, the sign of $x$ for our solution will not change, and therefore $|x(t)|$ eventually decreases monotonically.

        % \begin{figure}
        %     \centering
        %     \includegraphics[width=0.6\linewidth]{divided_region.png}
        %     \caption{The region $y'(t)>0$ separated by the line $18y+10x=0$.}
        %     \label{fig:y_prime_pos_sep_by_line}
        % \end{figure}

        Consequently, we may view $y$ as an invertible function of $x$ for sufficiently large values of $t$.
        Thus, considering $\frac{dy}{dx} = \frac{y'(t)}{x'(t)}$, we have by \eqref{eq:sys_xy} that 
        \[ 
            \frac{dy}{dx} = \frac{-6x^3 y - 18x^2 y^2 + 6x^2 + 18x y + 36y^2 - 6}{-6x\left(x^3 + 3x^2 y + 4x + 6y\right)}. 
        \]
        % Since $|x(t)|$ is bounded and eventually monotonically decreasing while $y(t)$ is unbounded and eventually monotonically increasing, we can consider the limit of $\frac{dy}{dx} \cdot \frac{1}{y}$ as $y\to \infty$, where we treat $x$ as a function of $y$.
        Using that $|x|$ is bounded, we get that
        \begin{equation}\label{eqn A}
             \lim_{y \to \infty}\frac{dy}{dx} \cdot \frac{1}{y} = \lim_{y \to \infty}\frac{x^2-2}{x\left(x^2+2\right)}. 
        \end{equation}
        Because $|x(t)|$ is bounded and monotonic, it must converge to some finite value as $t$ increases (equiv. $y\to \infty$).
        For the sake of contradiction, assume this limit is non-zero.
        % Thus as $y\to \infty$, 
        % This is equivalent to letting $y \to \infty$, and we see in this case that the above limit does exist in this case (let's call the limit $k$) since that would imply the expression above is (eventually) monotonic and bounded. 
        Then the limit \eqref{eqn A} converges to some finite value, and thus there exists a positive number $M$ such that, for sufficiently large values of $t$, we have
        \[
            \left | \frac{dy}{dx} \cdot \frac{1}{y} \right | < M.
        \]
        % Hence, the growth rate of $|\frac{dy}{dx}|$ is eventually dominated by a linear function.  
        Because $|x(t)|$ is bounded, the above inequality implies that $|y(t)|$ must also be bounded, which is a contradiction.
        Thus, we have $x\to 0$ as $t$ increases.
        %
        % \todo[inline]{Line separating new (above) from old (below).}
        %
        % Hence, we can assume $C$ is contained in the connected component of the region $y'>0$ above the line $18y+10x=0$. 
        % A computation in \textsc{MAPLE} shows that the region for which $y'>0$, $18y+10x>0$, and $y''<0$ is empty, meaning the connected component lying above this line is forward-invariant.
        % Furthermore, $x' > 0$ if $x < 0$ and $x' < 0$ if $x > 0$ in this component.  
        % Since $x=0$ is an invariant curve, note that our integral curve $C$ will always have either $x<0$ and $x'>0$ or $x>0$ and $x'<0$.  In other words, $|x(t)|$ eventually decreases monotonically and $y(t)$ eventually increases monotonically.  
        %
        %
        % Eventual monotonicity of $x(t)$ allows us to view $y$ as a function of $x$ for large values of $t$.  
        % We can thus consider the derivative $\frac{dy}{dx} = \frac{dy/dt}{dx/dt}$.
        % Observe that this quantity is always defined since $\frac{dx}{dt} \neq 0$.  
        % Furthermore, by \eqref{eq:sys_xy}, we have
        % \[ \frac{dy}{dx} = \frac{-6x^3 y - 18x^2 y^2 + 6x^2 + 18x y + 36y^2 - 6}{-6x\left(x^3 + 3x^2 y + 4x + 6y\right)}. \]
        %
        % Since $|x(t)|$ is bounded and eventually monotonically decreasing while $y(t)$ is unbounded and eventually monotonically increasing, we can consider the limit of $\frac{dy}{dx} \cdot \frac{1}{y}$ as $y\to \infty$, where we treat $x$ as a function of $y$.
        % Using that $|x|$ is bounded, we have
        % \begin{equation}\label{eqn A}
        %      \lim_{y \to \infty}\frac{dy}{dx} \cdot \frac{1}{y} = \lim_{y \to \infty}\frac{-18x^2+36}{-6x\left(3x^2+6\right)}. 
        % \end{equation}
        %
        %
        % Because $x(t)$ is bounded and eventually monotonic, it must converge to some finite value as $t \to \infty$ (equiv. $y\to \infty$).
        % For the sake of contradiction, assume this limit is non-zero.
        % % Thus as $y\to \infty$, 
        % % This is equivalent to letting $y \to \infty$, and we see in this case that the above limit does exist in this case (let's call the limit $k$) since that would imply the expression above is (eventually) monotonic and bounded. 
        % Then the limit \eqref{eqn A} converges to some $k$, and thus given $\epsilon > 0$, for sufficiently large values of $t$, we have that
        % \[
        %     \left | \frac{dy}{dx} \right | < |k|y+\epsilon
        % \]
        % Hence, the growth rate of $|\frac{dy}{dx}|$ is eventually dominated by a linear function.  
        % This means that as $y$ tends to infinity, $x$ must also tend to infinity. This means that if $|x(t)|$ is bounded, $|y(t)|$ must also be bounded, contradicting Lemma \ref{lem:compact}.
        % Therefore, $|x(t)|$ must go to zero monotonically as $t \to \infty$.
    \end{proof}
    
%Since $x=0$ is an integral curve, the uniqueness of solutions implies that $x$ remains positive along the flow. Consequently, we have $\frac{dP}{dt} > 0$. 

    We are now ready to prove the following case of Theorem \ref{mainthm:Ricpos}:

    \begin{theorem}\label{thm:P>0firstcase}
    Suppose $(x(t),y(t))$ is a solution of System \eqref{eq:sys_xy} whose initial condition lies in region $C$ or $F$ from \eqref{regions}.
    % \begin{enumerate}
    %     \item $x > -\sqrt{2}$, or \label{item:condition1}
    %     \item $y > -\displaystyle\frac{2x^2+3}{6x}$, $y > -\displaystyle\frac{3x^2+2}{8x}$, and $x<0$.\label{item:condition2}
    % \end{enumerate}
    Then the solution will satisfy $\Ric>0$ for all sufficiently large values of $t$.
\end{theorem}

\begin{proof}
    % First, if $3x^{2}+8xy+2<0$ and $2x^{2}+6xy+3>0$ (with $-\sqrt{6}<x<0$), or if $3x^{2}+8xy+2>0$ and $-\sqrt{2}<x<0$, then it can be verified with Maple's IsEmpty function that $P>0$.
    % Because these regions are bounded invariant curves, they are forward-invariant, and if the initial condition lies in this region, we have that $\Ric>0$ for all $t$.
    % Thus, we may assume that the initial condition satisfies either $x>0$, or $3x^{2}+8xy+2<0$ and $2x^{2}+6xy+3<0$ (with $x<0$).
    
    % We claim that $|x(t)|$ must be bounded.
    % If the initial condition satisfies $x>0$, then $x(t)>0$ for all $t$ because $x=0$ is an invariant curve.
    % Thus $x'(t) = -6x\left(x^3 + 3x^2 y + 4x + 6y\right)<0$ for all $t$, and hence $0<|x(t)|\leq x(0)$ for all $t$ in this case.
    % Suppose instead the initial condition satisfies $x<0$.
    % Then $x(t) < 0$ for all $t$, and one can check in Maple if $3x^{2}+8xy+2<0$ and $2x^{2}+6xy+3<0$, then $x'(t) > 0$.
    % Therefore, $x(0)\leq |x(t)|<0$ for all $t$ in this case, and hence $|x(t)|$ is bounded in both cases.
    
    % Now suppose for the sake of contradiction that that $(x(t),y(t))$ satisfies $P\leq 0$ for all $t$.
    By Lemma \ref{lemma D}, $|x(t)|$ converges monotonically to $0$ and $y(t)$ is unbounded and eventually grows monotonically as $t$ increases.
    %
%
    % Let $C=(x(t),y(t))$ be an integral curve that satisfies the conditions in the theorem. 
    % Assume, for contradiction, that $C$ never enters the $P>0$ region. 
    % So, from Lemma \ref{lem:compact}, we get that $C$ must escape any compact region. We will first prove that,
 %   
    % Claim 1: $C$ cannot escape to $x\rightarrow \pm \infty.$
%
    % \todo[inline]{Claim 1: $x$ is bounded. Proof of Claim 1: If $x$ is unbounded from below, then eventually $x <-\sqrt{6}$. However, if $x <-\sqrt{6}$ and $y>-\frac{3x^2+2}{8x}$ (which it must since $C$ cannot cross an invariant curve), then $dx/dt > 0$ (Checked with MAPLE). Thus $x$ is bounded below by its initial value. $x$ must be bounded from above because we start with $x<0$ and the curve $x=0$ is an invariant curve for the system. On the other hand, if we DID start with $x>0$, then since $y>0$, it follows that $dxdt<0$ by inspection, and hence x is bounded above by its initial condition.}
%
  %  
 %   
    %  implies that as $t$ goes to infinity, $x(t)$ goes to $0$. By comparing the boundary of the region $P$ and behavior of $C,$ we will prove that
%
    % Claim 2: As $x \rightarrow 0,$  $C$ enters and stays in the region $P>0.$
 %  
%
% Let us first prove claim 1. We start with the case when $x$ approaches $-\infty.$ Note that, by our hypothesis, this forces that $C$ lies above the invariant curve. Now, we consider the limit of the invariant curve $y=-\frac{3x^2+2}{8x}$ as $x\to-\infty$. In this limit, the invariant curve approaches the line defined by $8y+3x=0$. Since our integral curve $C$ can never cross the invariant curve $y=-\frac{3x^2+2}{8x}$ and the curve $y=-\frac{3x^2+2}{8x}$ lies above $8y+3x=0$, $C$ cannot intersect the line $8y+3x=0$ as well, hence, in order to avoid intersecting with this line, as $x$ approaches $-\infty$, then $y$ coordinate of $C$ must approach $\infty$.\todo[inline]{maybe graphs of these curves?} That is, the escape direction “down-left” is impossible for $C$; So $C$ must go “up-left”. We examine $y'(t)$ to show that this cannot happen. From \eqref{eq:sys_xy}, we find that the roots of $y'(t)=0$ are:
% \begin{align*}
% y_1 &= \frac{x^3 - 3x - \sqrt{x^6 + 6x^4 - 27x^2 + 24}}{6(2-x^2)}, \\
% y_2 &= \frac{x^3 - 3x + \sqrt{x^6 + 6x^4 - 27x^2 + 24}}{6(2-x^2)}
% \end{align*}
% Computing 
% \begin{align*}
%     \lim_{x \to -\infty} y_1=\lim_{x \to -\infty} \frac{x^3 - 3x - |x^3|\sqrt{1 + \frac{6}{x^2} - \frac{27}{x^4} + \frac{24}{x^6}}}{6(2-x^2)}= 
% \lim_{x \to -\infty} \frac{2x^3-3x}{12-6x^2}=-\frac{x}{3} \\[1mm]
% \lim_{x \to -\infty} y_2=\lim_{x \to -\infty} \frac{x^3 - 3x +|x^3|\sqrt{1 + \frac{6}{x^2} - \frac{27}{x^4} + \frac{24}{x^6}}}{6(2-x^2)}= 
% \lim_{x \to -\infty} \frac{-3x}{12-6x^2}=-\frac{2}{x} \\[1mm]
% \end{align*}
% From \eqref{eq:sys_xy}, observe that as $x \to -\infty$ we have
% \[
% y' = A (y - y_1)(y - y_2)
% \]
% for some constant $A<0$. In particular, $y'>0$ only when $y$ lies between the curves $y_1$ and $y_2$. For large negative values of $x$, the limit computation above shows that
% \[
% y_1 \longrightarrow \frac{-x}{3} 
% \qquad \text{and} \qquad
% y_2 \longrightarrow 0.
% \]
% Thus both $y_1$ and $y_2$ lie strictly below the asymptote of the invariant curve,
% \[
% 8y + 3x = 0.
% \]
% Consequently, for $x$ sufficiently negative, the entire region where $y'>0$ lies below the invariant curve. This contradicts our assumption that the integral curve $C$ lies above the invariant curve. Therefore, the case $x \to -\infty$ is impossible.
%
%
% Next, we assume $x \to \infty.$ From \eqref{eq:sys_xy}, \[x'(t) = -6x\left(x^3 + 3x^2 y + 4x + 6y\right).\]
% The terms inside the parenthesis will be positive since $x$ and $y$ are constrained to be positive, resulting in $x'(t)<0$. Therefore, $x$ is decreasing and hence cannot go to $\infty,$ which is a contradiction.
%
% Suppose $C$ lies below the invariant curve with $x>-\sqrt{2}.$ The invariant curve $y=-\frac{3x^2+2}{8x}$ tangentially intersects the admissible branch of the bounding hyperbola $2y(3y+x)=1$ at $\left(-\sqrt{2},\frac{1}{\sqrt{2}}\right).$ Since an integral curve cannot cross either the invariant curve or the bounding hyperbola, the curve $C$ must remain in the region between these two curves. 
% % On the relevant branch of the hyperbola, we have $y>0.$ Because $C$ lies above this hyperbola, it follows that $C$ also satisfies $y>0.$ 
% % Consequently, the invariant curve evaluated along $C$ must also satisfy $y>0.$ 
% From the equation of the invariant curve, since $y>0$, we have $x<0.$ Therefore, the portion of $C$ under consideration is confined to the interval $-\sqrt{2}<x <0.$ This completes the proof of claim 1.
%
% Now, it remains to prove claim 2.
%
Thus, we may again treat $y$ as a function of $x$, and we may consider
% \[
% \displaystyle\lim_{x\rightarrow 0}\frac{dy}{dx}\cdot \frac{x}{y}
% = \displaystyle\lim_{x\rightarrow 0}\frac{y'(t)}{x'(t)}\cdot \frac{x}{y}
% = \displaystyle\lim_{x\rightarrow 0}\frac{(-6x^{3}y - 18x^{2}y^{\,2} + 6x^{2} + 18xy + 36y^{2} - 6)x}
% {-6xy\,(x^{3} + 3x^{2}y + 4x + 6y)}\]
% \[ = \lim_{x\rightarrow 0}\frac{y^2(36x-18x^3)+...}{-6xy^2(3x^2+6)+...}
% =\displaystyle\lim_{x\rightarrow 0} \frac{36x-18x^3}{-6x(3x^2+6)} = -1.
%\frac{36y^2 - 6}{-36 xy}< \frac{36y^2}{-36 xy}=-\frac{y}{x}.
% \]
    \begin{align*}
        \frac{dy}{dx}\cdot \frac{x}{y} 
        = \frac{y'(t)}{x'(t)}\cdot \frac{x}{y} 
        &= \frac{ (3x^2-6)y^2 + (x^3 - 3x)y + 1 - x^2}{ (3x^2+6)y^2 + (x^3  + 4x)y }.
        %&= \frac{ 3x^2-6 + (x^3 - 3x)/y + (1-x^2)/y^2}{ 3x^2+6 + (x^3  + 4x)/y }.
    \end{align*}
    Because $y\to\infty$ as $x\to 0$ (equiv. as $t$ increases), we get that
    \[
        \lim_{x \to 0} \left|\frac{dy}{dx}\cdot \frac{x}{y}\right| \leq  1.
    \]
    Hence given $\epsilon>0$, for small enough values of $x$, we have
    \[
        \left|\frac{dy}{dx}\right| < (1+\epsilon)\frac{y}{|x|}
    \]
    % {\color{red}Because $\frac{dy}{dx}>0$ when $x<0$ and $\frac{dy}{dx}<0$ when $x>0$, it follows that
    % \[
    %     \begin{cases}
    %         \frac{dy}{dx} < -(1+\epsilon)\frac{y}{x} & \text{if }x<0,\\
    %         \frac{dy}{dx} > -(1+\epsilon)\frac{y}{x} & \text{if }x>0.
    %     \end{cases}
    % \]}
    Now, solutions of the differential equation $\frac{dy}{dx} = -(1+\epsilon)\frac{y}{x}$ are given by $y = A|x|^{-(1+\epsilon)}$ for some constant $A$.
    Thus, for sufficiently large values of $t$, the solution $(x(t),y(t))$ satisfies $y(x) < A|x|^{-(1+\epsilon)}$.
    In particular, choosing $\epsilon <1$, we have $A|x|^{-(1+\epsilon)}<y_2(x)$ for sufficiently small values of $x$, where $y_2$ is the curve described in Remark \ref{rmk 1}.
    Thus, for large enough values of $t$, we get that $(x(t),y(t))$ satisfies $y<y_3<y_2$, and hence $P>0$, which implies $\Ric>0$.
% {\color{red}
%     Wait... Does this even make sense for $x<0$?
%     I suppose if $1+\epsilon$ is rational with odd denominator, then $x^{-(1+\epsilon)}$ is defined on the set of negative real.
% }
%
% But
% \[
% \frac{dy}{dx} = -\,\frac{y}{x}
% \]
% has solution
% \[
% y(x) = \frac{B}{x},
% \]
% for some positive constant $B.$ So, the curve $y(x) = \frac{B}{x}$ serves as an upper bound for the growth of $y$ along the actual
% integral curve.
%
% On the other hand, from Remark \ref{rmk 1}, the region $P>0$ is bounded above the curve
% \[
% y_2(x)
% = -\frac{x}{3}
%   - \frac{1}{2x}
%   + \frac{\sqrt{3}}{2x^{2}}.
% \]
% As $x\to 0$, the dominant term in $y_2(x)$ is $\frac{\sqrt{3}}{2x^{2}}$, so
% \[
% y_2(x) \sim \frac{\sqrt{3}}{2x^{2}}.
% \]
% Since $x^{-2}>x^{-1}$ for $x$ near zero,
% \[
% y_2(x) > \frac{B}{x}
% \qquad \text{for all sufficiently small } |x|.
% \]
% Consequently, once $|x|$ is small enough, the integral curve $C$ lies strictly
% below the boundary curve $y_2(x)$, and therefore $C$ remains entirely in the region
% $P>0$.  This proves claim 2.
\end{proof}

% \todo[inline]{This is barely a proof by contradiction. See my note next to the previous lemma. I think there is a way to make these arguments more cohesive.}

It remains to show Theorem \ref{mainthm:Ricpos} holds in region $A$.
First, we show the following:

\begin{lemma}\label{lemma E}
    Suppose $(x(t),y(t))$ is a solution of System \eqref{eq:sys_xy} whose initial condition lies in region $A$ from \eqref{regions}.
    Then $x(t)$ decreases monotonically for all sufficiently large values of $t$, and $x(t) \to -\infty$ as $t$ increases.
\end{lemma}

\begin{proof}
    % If we get into the $P\ge 0$ region, it can be verified with maple that $x'(t)<0$ and $y'(t)>0$. Therefore, $x$ monotonically decreases. Since $P\ge 0$ is a forward-invariant region with boundary $P=0$ given by the curve 
    % $$
    % -\frac{x}{3}-\frac{1}{2x}+\frac{\sqrt{3}}{2x^2},
    % $$
    % we see that $y$ cannot go to $\infty$ unless $x$ tends to $-\infty$. Since we exclude the possibility of the integral curve being contained in a compact region, we reach the conclusion that $x$ must tend to $-\infty$ monotonically. Therefore, suppose that we stay in the $P<0$ region.
    % 
    By Remark \ref{rem:compact}, we know that our solution is unbounded. 
    Within the region $A$, because it is bounded by the invariant curves $3x^{2}+8xy+2 = 0$ and $2x^2 + 6xy + 3 = 0$ (with $x<-\sqrt{6}$), a solution contained in $A$ is unbounded only if $x(t)\to -\infty$ and $y(t) \to \infty$ as $t$ increases.
    In particular, there exists $t$ for which $x'(t) < 0$ and $y'(t) > 0$.
    Furthermore, it is verifiable with Maple's IsEmpty function that in region $A$, if $y'(t)\geq 0$, then $x'(t)<0$ and $x''(t)<0$.
    Thus the subset of $A$ for which $x'(t)<0$ is forward-invariant.
    %
    %   Curve $B$ intersects the bounding hyperbola $2y(3y+x)=1$ at $\left(-\sqrt{2},\frac{\sqrt{2}}{2}\right)$. Since no integral curve can cross the invariant curve nor the bounding hyperbola, and all initial conditions are taken to lie above this boundary hyperbola, $x$ is bounded above by $-\sqrt{2}$. The fact that all integral curves remain bounded away from $x=0$, where $B$ has a vertical asymptote, and cannot cross the invariant curve, implies that $y$ is bounded as a function of $x$. In particular, $y$ cannot approach $\infty$ unless $x \to -\infty$. Once again, since no integral curve can cross the bounding hyperbola, its $x$-component cannot approach $-\infty$ with $y$ bounded. By hypothesis, the integral curve does not stay in a compact region. This leaves us with no other option but the one where $y\to\infty$ while $x\to-\infty$. 
    %   Therefore, there exists $t$ for which $x'(t)<0$ and $y'(t)>0$.
    %   In the next paragraph, we prove eventual monotonicity of the $x$-component.
      %
    %  Suppose first that we start at $x\in[-\sqrt{6},-\sqrt{2}).$ It can be seen that, when $-\sqrt{6}<x<-\sqrt{2}$, curve $B$ and every point under it stays inside of the region where $x'<0$ and $y'>0$, and then curves $x'=0, y'=0$ and $B$ intersect at $x=-\sqrt{6}$.
      %
    %   \begin{figure}[htbp]
    %         \centering
    %         \includegraphics[width=0.6\textwidth]{betweenintersections.png}
    %         \caption{The red curve $B$ is inside of the colored regions, where the purple region shows where $x'<0$, while the green region shows where $y'>0$. The intersection points of curves $B$, $x'=0$ and $y'=0$ is on the left, while the $\left(-\sqrt{2}, \frac{\sqrt{2}}{2}\right)$ fixed point is on the right).}
    %         \label{fig:betweenintersections}
    % \end{figure}
%
      %
    %   Since no integral curve from under $B$ can cross this invariant curve, by continuity of the integral curve and openness of the $x'<0$ and $y'>0$ regions, we see that any integral curve must reach points where $x'<0$, $y'>0$, and $x<-\sqrt{6}$. One can check with MAPLE that, when $x<-\sqrt{6}$, the region where $y'>0$ is forward-invariant and contained in the region where $x'<0$ (see Figure \ref{fig:purplegreen}).
%
      %
    %   \begin{figure}[htbp]
    %         \centering
    %         \includegraphics[width=0.6\textwidth]{purplegreen.png}
    %         \caption{The purple region shows where $x'<0$, while the green region shows where $y'>0$. One can see that the green region lies inside of the purple one.}
    %         \label{fig:purplegreen}
    % \end{figure}
%
    %
    %   Thus, from the point when the integral curve crosses the $x=-\sqrt{6}$ line, $x$ will monotonically decrease, while $y$ will monotonically increase. 
%
    %   Now, suppose that we start at $x<-\sqrt{6}$. If we start in the $y'<0, x'<0$ region (purple region $\backslash$ green region), then the slope of the integral curve is positive. Meanwhile, Figure \ref{fig:purplegreen} shows us that the slope of the $y'=0$ curve is negative, and in fact one could check that it will approach $-\frac{1}{3}$. 
    %   Therefore, the integral curve will eventually reach the $y'>0$ region (green), which would be forward-invariant and contained in the $x'<0$ region for $x<-\sqrt{6}$. 
      %
    %   % On the other hand, if we start in the white region ($x'>0, y'<0$) and assume that we stay in it, then the $x$-component and $y$-component of the integral curve will increase and decrease, respectively. Since the invariant curve $B$ intersects the $x'=0$ and $y'=0$ curves at $x=-\sqrt{6}$, the integral curve will either enter (see Figure \ref{fig:betweenintersections}) the $x'<0$ region or approach some kind of limit point. However, the latter would violate Lemma \ref{lem:compact}. 
      %
    %   % Therefore, the integral curve must enter the $x'<0$ region and, by the previous argument, it will eventually enter the $y'>0$ region, which is forward-invariant and contained in the $x'<0$ region for $x<-\sqrt{6}$. 
    %   Thus, there will exist a $T<\infty$, such that for all $t\geq T$, the $x$-component and $y$-component will monotonically decrease and increase, respectively, along the integral curve. This proves the statement.
\end{proof}


    % Finally, we prove the following, which implies Theorem \ref{mainthm:Ricpos}:

    % \begin{theorem}
    %     For every solution $(a(t),b(t))$ to the System \eqref{eq:sys_ab}, there exists $t_0$ such that the metric $g(t)$ has $\Ric > 0$ for all $t \geq t_0$.
    % \end{theorem}

    Finally, we prove the last case of Theorem \ref{mainthm:Ricpos}:
    
    \begin{theorem}\label{thm:enter-P-positive-xy}
        Suppose $(x(t),y(t))$ is a solution of System \eqref{eq:sys_xy} whose initial condition lies in region $A$ from \eqref{regions}.
        Then the solution will satisfy $\Ric>0$ for all sufficiently large values of $t$.
    \end{theorem}

    \begin{proof}
        By Lemma \ref{lemma E}, $x(t)$ eventually decreases monotonically, so we may regard $y(t)$ and $P(t)$ as functions of $x$.
        From \eqref{eq:sys_xy} and \eqref{eq:Pxy}, we get that
        \begin{equation}\label{eq:dP_dx_exact}
            \begin{aligned}
                \frac{dP}{dx}&=
                -\frac{72x^{2}\big(6xy+2x^{2}+3\big)^{2}\big(x^{3}+3x^{2}y+2y+2x\big)}
                {6x\big(x^{3}+3x^{2}y+4x+6y\big)}\\
                &= -12x\,(2x^2+6xy+3)^2
                \left(1-\frac{2x+4y}{x^3+3x^2y+4x+6y}\right).
            \end{aligned}
        \end{equation}
        Suppose the solution initially satisfies $P\leq 0$.
        Defining 
        \[ 
            P_1 = 2x^2+6xy+3 \quad \text{ and } \quad P_2=1-\frac{2x+4y}{x^3+3x^2y+4x+6y},
        \]
        we will first estimate $P_1$ and $P_2$ in order to bound $\frac{dP}{dx}=-12x(P_1)^2 P_2$ above by a negative function of $x$, which corresponds to a positive lower bound on $\frac{dP}{dt}$ since $x(t)$ decreases monotonically.
        %we will use this bound to justify that $P(t)>0$ for sufficiently large values of $t$.

        \begin{figure}[H]
            \centering
            \includegraphics[width=0.45\textwidth]{P_nonpositive.eps}
            \caption{The region where $P\leq0$ is bounded by $y=y_2(x)$.}
            \label{fig:negativep}
        \end{figure}
        
        First notice that in the region where $P\leq 0$, for a fixed value of $x$, the minimal value of $y$ occurs on the curve where $P=0$, which is given by $y_2(x) = -\frac{x}{3} - \frac{1}{2x} + \frac{\sqrt{3}}{2x^2}$ in Remark \ref{rmk 1}; see Figure \ref{fig:negativep}.
        We also have that $\frac{\partial P_1}{\partial y}=6x<0$ in the region $A$.
        Thus for a fixed value of $x$, the maximal value of $P_1$ is attained on $y_2(x)$.
        Hence for all $x$,
        \begin{equation}\label{eq:P1}
            P_1 = 2x^2 + 6xy + 3 \leq \frac{3\sqrt{3}}{x} < 0.
        \end{equation}

        Next notice, because $x<-\sqrt{6}$ in region $A$, we have
        \[
            \frac{\partial P_2}{\partial y} = \frac{2x(x^2-2)}{(x^3+3x^2y+4x+6y)^2} < 0.
        \]
        So similarly, for a fixed value of $x$, the maximal value of $P_2$ is attained on $y_2(x)$.
        Hence for all $x$,
        \[
            P_2 = 1-\frac{2x+4y}{x^3+3x^2y+4x+6y} \leq - \frac{x^3-9\sqrt{3}x^2+6x-6\sqrt{3}}{3x^3+9\sqrt{3}x^2-18x+18\sqrt{3}}.
        \]
        Consequently, 
        \begin{equation}\label{eq:P2}
            \lim_{x\to-\infty} P_2 \leq -\frac{1}{3}.
        \end{equation}

        Combining \eqref{eq:P1} and \eqref{eq:P2}, for any sufficiently small $\epsilon>0$ and for sufficiently large values of $x$, we have 
        \[
            \frac{dP}{dx} \leq -12x \of{\tfrac{3\sqrt{3}}{x}}^2 \of{-\tfrac{1}{3}+\epsilon} = \frac{324 (1-3\epsilon)}{3x} < 0.
        \]
        Consequently, there exist constants $C_1, C_2$ such that $C_1>0$ and for sufficiently large values of $x$,
        \[P > C_1\ln |x|+C_2.\]
        Because $x\to -\infty$ as $t$ increases, we conclude that $P > 0$ for sufficiently large values of $t$.
        %
%         We argue by contradiction.  Suppose the trajectory never enters the region $P>0$, i.e.\ it stays in the region $\{P\le 0\}$ for all forward time.  
%         By Lemma \ref{lemma E}, under this hypothesis \(x(t)\) is monotone and satisfies
%         \[
%         x(t)\longrightarrow -\infty \qquad\text{as }t\to+\infty.
%         \]
%         Monotonicity of \(x\) allows us to regard \(y\) and \(P\) as functions of \(x\) along the trajectory.
%         Combining \eqref{eq:sys_xy} and \eqref{eq:Pxy}, we get that
%         % \begin{equation}%\label{eq:dP_dx_exact}
%         % \begin{aligned}
%         % \frac{dP}{dx}&=
%         % -\frac{36x^{2}\big(6xy+2x^{2}+3\big)^{2}\big(x^{3}+3x^{2}y+2y+2x\big)}
%         % {6x\big(x^{3}+3x^{2}y+4x+6y\big)}\\&= -6x\,(6yx+2x^2+3)^2
%         % \left(1-\frac{2x+4y}{x^3+3x^2y+4x+6y}\right).
%         % \end{aligned}
%         % \end{equation}
% %
%         % {\color{blue}
%             \begin{equation}%\label{eq:dP_dx_exact}
%                 \begin{aligned}
%                     \frac{dP}{dx}&=
%                     -\frac{72x^{2}\big(6xy+2x^{2}+3\big)^{2}\big(x^{3}+3x^{2}y+2y+2x\big)}
%                     {6x\big(x^{3}+3x^{2}y+4x+6y\big)}\\
%                     &= -12x\,(6xy+2x^2+3)^2
%                     \left(1-\frac{2x+4y}{x^3+3x^2y+4x+6y}\right).
%                 \end{aligned}
%             \end{equation}
%         % }
        %
%         We now provide an upper bound for \eqref{eq:dP_dx_exact}.
%         Because $x<0$, it is sufficient to find an upper bound for $6xy+2x^2+3$ and $1-\frac{2x+4y}{x^3+3x^2y+4x+6y}$.
%         % the right-hand side of equation \eqref{eq:dP_dx_exact}. 
        %
        %
        %
%
        %
%     First, we obtain an upper bound on the limit of $1-\frac{2x+4y}{x^3+3x^2y+4x+6y}$ as $x\to -\infty$.
%     Now, since 
%         $$
%         \frac{\partial }{\partial y}\left(1-\frac{2x+4y}{x^3+3x^2y+4x+6y}\right)=\frac{2x(x^2-2)}{(x^3+3x^2y+4x+6y)^2}<0,
%         $$
%     we see that minimizing $y$ with fixed $x$ actually maximizes the value of $1-\frac{2x+4y}{x^3+3x^2y+4x+6y}$.
%     Because $C$ stays in the region \(P\le 0\), it remains above the curve where \(P=0\), so the maximum of $1-\frac{2x+4y}{x^3+3x^2y+4x+6y}$ is attained on the $P=0$ curve, where $y$ has the lowest values possible. 
%     On $P=0,$ Remark \ref{rmk 1} implies that
%         \[
%         y = -\frac{x}{3} - \frac{1}{2x} + \frac{\sqrt{3}}{2x^2}.
%         \]  
%         Substituting this into $ 1-\frac{2x+4y}{x^3+3x^2y+4x+6y}$ and simplifying yields
        %
%         \[ 1-\frac{2x+4y}{x^3+3x^2y+4x+6y} \leq 1-\frac{\frac{4x}{3}-\frac{2}{x}+\frac{2\sqrt 3}{x^2}}{\frac{x}{2}+\frac{\sqrt 3}{2}-\frac{3}{x}+\frac{3\sqrt 3}{x^2}}. \]
        %
%         Taking $x\to-\infty$, we obtain the upper bound
   %     
%         \begin{equation}\label{eq:limit_upper_bound}
%             \lim_{x \to -\infty} \left ( 1-\frac{2x+4y}{x^3+3x^2y+4x+6y} \right ) \leq -\frac{5}{3}.
%         \end{equation} 
        %
    %
%         Next, we obtain an upper bound for $(6xy+2x^{2}+3)^2$ in terms of $x$.
%         Define \(\Delta(x,y) =6xy+2x^{2}+3\).
%         Notice that, $\frac{\partial \Delta}{\partial y}=6x<0$. 
%         Thus for a fixed value of $x$, the maximum of $\Delta$ is attained on the curve $P=0$, where $y=-\frac{x}{3} - \frac{1}{2x} + \frac{\sqrt{3}}{2x^2}$. 
%         % \[
%         %     6 \left(-\frac{x}{3} - \frac{1}{2x} + \frac{\sqrt{3}}{2x^2} \right) x + 2x^2 + 3=-2x^2 - 3 + \frac{3\sqrt{3}}{x} + 2x^2 + 3.
%         % \]
%         %Expanding term by term:
%         %\[
%         %6 \left(-\frac{x}{3} \cdot x \right) = -2x^2, \quad
%         %6 \left(-\frac{1}{2x} \cdot x \right) = -3, \quad
%         %6 \left(\frac{\sqrt{3}}{2x^2} \cdot x \right) = \frac{3\sqrt{3}}{x}.
%         %\]
%         Hence, 
%         \[
%             6xy + 2x^2 + 3 \leq  \frac{3\sqrt{3}}{x}<0,
%         \]
%         % which provides an upper bound for $\Delta(x,y)$. Since the upper bound is negative, so is the lower bound. Thus, when squaring, due to the sign change, $(\frac{3\sqrt{3}}{x})^2=\frac{27}{x^2}$ will become the lower bound on $[\Delta(x,y)]^2$.
%         which implies that 
%         \begin{equation}\label{eq:second_bound}
%             (6xy + 2x^2 + 3)^2 \geq  \frac{27}{x^2}.
%         \end{equation}
            %
%         Combining \eqref{eq:limit_upper_bound} and \eqref{eq:second_bound}, for any sufficiently small $\epsilon > 0$, there exists a number $N < 0$ such that for any $x < N$, we have
        %
%         % \[ \frac{dP}{dx} < -6x \left ( \frac{3\sqrt{3}}{x}\right)^2 \left (\frac{-5}{3} + \epsilon \right ) = \frac{-162}{x} \left (\frac{-5}{3} + \epsilon \right )<0  \]
%
%         % {\color{blue} 
%             \[ \frac{dP}{dx} < -12x \left ( \frac{3\sqrt{3}}{x}\right)^2 \left (\frac{-5}{3} + \epsilon \right ) = \frac{-324}{x} \left (\frac{-5}{3} + \epsilon \right )<0.  \]
%         % }
    %
%         % Therefore, 
%         % $$
%         % \frac{dP}{d(-x)}>\frac{162}{x} \left (\frac{-5}{3} + \epsilon \right )>0
%         % $$
        %
%         Consequently, there exist constants $C_1, C_2$ 
%         % {\color{blue} 
%             such that $C_1>0$ 
%         % }
%         and for $x < N$,
        %
%         \[P > C_1\ln |x|+C_2\]
        %
%         which implies $P > 0$ as $x \to -\infty$.
    \end{proof}

\bibliographystyle{abbrv}
\bibliography{biblio}





\end{document}
